%%%%%%%%%%%%%%%%%%%%%%%%%%%%%%%%%%%%%%%%%%%%%%%%%%%%%%%%%%%%%%%%%%%%%%%%%
%% PLANTILLA LATEX PARA TRABAJO FIN DE GRADO
%% Universidad Camilo José Cela
%% Facultad de Tecnología y Emprendimiento
%% Grado en Ingeniería Informática Biomédica
%%%%%%%%%%%%%%%%%%%%%%%%%%%%%%%%%%%%%%%%%%%%%%%%%%%%%%%%%%%%%%%%%%%%%%%%%

\documentclass[12pt, a4paper, oneside]{report}

%% ==================== PAQUETES ==================== %%
\usepackage[utf8]{inputenc}
\usepackage[T1]{fontenc}
\usepackage[spanish, es-tabla]{babel}
\usepackage{mathptmx}                    % Fuente Times New Roman
\usepackage{geometry}
\usepackage{setspace}
\usepackage{titlesec}
\usepackage{titletoc}
\usepackage{fancyhdr}
\usepackage{graphicx}
\usepackage{xcolor}
\usepackage{hyperref}
\usepackage{bookmark}
\usepackage{array}
\usepackage{booktabs}
\usepackage{longtable}
\usepackage{multirow}
\usepackage{float}
\usepackage{caption}
\usepackage{subcaption}
\usepackage{listings}
\usepackage{amsmath, amssymb, amsfonts}
\usepackage{enumitem}
\usepackage{tocloft}
\usepackage{appendix}
\usepackage{pdfpages}
\usepackage{csquotes}
\usepackage[backend=biber, style=apa, sorting=nyt]{biblatex}
\DeclareLanguageMapping{spanish}{spanish-apa}
\usepackage{etoolbox}
\usepackage{ragged2e}
\usepackage{microtype}
\usepackage{lastpage}

%% ==================== DATOS DEL TFG ==================== %%
\newcommand{\tituloes}{DESARROLLO DE UN MODELO DE INTELIGENCIA ARTIFICIAL PARA LA PREDICCIÓN TEMPRANA DE LESIÓN RENAL AGUDA UTILIZANDO DATOS CLÍNICOS HOSPITALARIOS}
\newcommand{\tituloen}{DEVELOPMENT OF AN ARTIFICIAL INTELLIGENCE MODEL FOR EARLY PREDICTION OF ACUTE KIDNEY INJURY USING HOSPITAL CLINICAL DATA}
\newcommand{\alumno}{Isabella Cristancho}
\newcommand{\emailalumno}{isabella.cristancho@alumno.ucjc.edu}
\newcommand{\grado}{Grado en Ingeniería Informática Biomédica}
\newcommand{\cursoacademico}{Curso académico 2022-2026}
\newcommand{\tutor}{Xavier Tarrago}
\newcommand{\facultad}{Facultad de Tecnología y Emprendimiento}
\newcommand{\universidad}{Universidad Camilo José Cela}

%% Ruta al logo (ajustar según tu estructura de carpetas)
\newcommand{\logopath}{img/logoUCJC.png}

%% ==================== CONFIGURACIÓN DE PÁGINA ==================== %%
\geometry{
    a4paper,
    left=3cm,
    right=2.5cm,
    top=2cm,
    bottom=2.5cm,
    headheight=28pt,
    headsep=10pt,
    footskip=30pt
}

%% Interlineado 1.5
\onehalfspacing

%% ==================== COLORES INSTITUCIONALES ==================== %%
\definecolor{ucjcgranate}{RGB}{128, 0, 64}    % Color granate UCJC
\definecolor{ucjcgray}{RGB}{128, 128, 128}
\definecolor{linkblue}{RGB}{0, 0, 128}

%% ==================== CONFIGURACIÓN DE HYPERREF ==================== %%
\hypersetup{
    colorlinks=true,
    linkcolor=black,
    citecolor=ucjcgranate,
    urlcolor=linkblue,
    bookmarksnumbered=true,
    bookmarksopen=true,
    pdfauthor={\alumno},
    pdftitle={\tituloes},
    pdfsubject={Trabajo Fin de Grado},
    pdfkeywords={Inteligencia Artificial, Lesión Renal Aguda, Machine Learning, Predicción Clínica}
}

%% ==================== CONFIGURACIÓN DE TÍTULOS ==================== %%
\titleformat{\chapter}[hang]
    {\normalfont\bfseries\LARGE}
    {\thechapter.}
    {0.5em}
    {}
\titlespacing*{\chapter}{0pt}{0pt}{15pt}

\titleformat{\section}
    {\normalfont\bfseries\large}
    {\thesection}
    {1em}
    {}
\titlespacing*{\section}{0pt}{10pt}{5pt}

\titleformat{\subsection}
    {\normalfont\bfseries\normalsize}
    {\thesubsection}
    {1em}
    {}
\titlespacing*{\subsection}{0pt}{10pt}{4pt}

\titleformat{\subsubsection}
    {\normalfont\itshape\normalsize}
    {\thesubsubsection}
    {1em}
    {}

%% ==================== ENCABEZADOS Y PIES DE PÁGINA ==================== %%

%% Estilo principal para páginas del documento
\fancypagestyle{tfgstyle}{
    \fancyhf{} % Limpiar todo
    %% ENCABEZADO
    % Logo a la izquierda
    \fancyhead[L]{\includegraphics[height=1cm]{\logopath}}
    % Texto centrado
    \fancyhead[C]{\textbf{TRABAJO FIN DE GRADO}}
    % Nada a la derecha del encabezado
    \fancyhead[R]{}
    % Línea debajo del encabezado
    \renewcommand{\headrulewidth}{1pt}
    %% PIE DE PÁGINA
    % Texto a la izquierda
    \fancyfoot[L]{TRABAJO FIN DE GRADO}
    % Número de página a la derecha
    \fancyfoot[R]{\thepage}
    % Línea encima del pie
    \renewcommand{\footrulewidth}{1pt}
}

%% Estilo para primera página de capítulos (plain)
\fancypagestyle{plain}{
    \fancyhf{}
    %% ENCABEZADO
    \fancyhead[L]{\includegraphics[height=1cm]{\logopath}}
    \fancyhead[C]{\textbf{TRABAJO FIN DE GRADO}}
    \fancyhead[R]{}
    \renewcommand{\headrulewidth}{1pt}
    %% PIE DE PÁGINA
    \fancyfoot[L]{TRABAJO FIN DE GRADO}
    \fancyfoot[R]{\thepage}
    \renewcommand{\footrulewidth}{1pt}
}

%% Estilo vacío para portada y páginas especiales
\fancypagestyle{empty}{
    \fancyhf{}
    \renewcommand{\headrulewidth}{0pt}
    \renewcommand{\footrulewidth}{0pt}
}

%% Aplicar estilo principal por defecto
\pagestyle{tfgstyle}

%% ==================== CONFIGURACIÓN DE ÍNDICES ==================== %%
\renewcommand{\cftchapfont}{\bfseries}
\renewcommand{\cftsecfont}{\normalfont}
\renewcommand{\cftsubsecfont}{\normalfont}
\renewcommand{\cftchapleader}{\cftdotfill{\cftdotsep}}

% Reduce el espacio enorme antes y después del título de los índices
\setlength{\cftbeforetoctitleskip}{0pt}
\setlength{\cftaftertoctitleskip}{12pt}
\setlength{\cftbeforeloftitleskip}{0pt}
\setlength{\cftafterloftitleskip}{12pt}
\setlength{\cftbeforelottitleskip}{0pt}
\setlength{\cftafterlottitleskip}{12pt}

%% ==================== CONFIGURACIÓN DE FIGURAS Y TABLAS ==================== %%
\captionsetup{
    font=small,
    labelfont=bf,
    format=hang,
    justification=justified
}

\captionsetup[table]{position=above}
%\captionsetup[figure]{position=below}

%% ==================== CONFIGURACIÓN DE CÓDIGO ==================== %%
\lstset{
    basicstyle=\ttfamily\footnotesize,
    breaklines=true,
    frame=single,
    numbers=left,
    numberstyle=\tiny\color{ucjcgray},
    keywordstyle=\color{ucjcgranate}\bfseries,
    commentstyle=\color{ucjcgray}\itshape,
    stringstyle=\color{red!60!black},
    showstringspaces=false,
    tabsize=4,
    captionpos=b,
    xleftmargin=2em,
    framexleftmargin=1.5em
}

%% ==================== CONFIGURACIÓN DE LISTAS ==================== %%
\setlist[itemize]{nosep, leftmargin=*}
\setlist[enumerate]{nosep, leftmargin=*}

%% ==================== BIBLIOGRAFÍA ==================== %%
\addbibresource{referencias.bib}  % Archivo de bibliografía

%% ==================== COMANDOS PERSONALIZADOS ==================== %%
\newcommand{\nota}[1]{\textcolor{red}{\textbf{[NOTA: #1]}}}

%%%%%%%%%%%%%%%%%%%%%%%%%%%%%%%%%%%%%%%%%%%%%%%%%%%%%%%%%%%%%%%%%%%%%%%%%
%% INICIO DEL DOCUMENTO
%%%%%%%%%%%%%%%%%%%%%%%%%%%%%%%%%%%%%%%%%%%%%%%%%%%%%%%%%%%%%%%%%%%%%%%%%

\begin{document}

%% ==================== PORTADA ==================== %%
\begin{titlepage}
    \thispagestyle{empty}
    \begin{center}
        
        %% Logo de la Universidad
        \includegraphics[width=8cm]{\logopath}
        
        \vspace{1cm}
        
        {\Large\bfseries TRABAJO FIN DE GRADO}
        
        \vspace{0.5cm}
        
        \rule{\textwidth}{1.5pt}
        
        \vspace{1cm}
        
        {\large\MakeUppercase{\grado}}
        
        \vspace{1.5cm}
        
        {\Large\bfseries TRABAJO FIN DE GRADO}
        
        \vspace{1cm}
        
        {\LARGE\bfseries\color{ucjcgranate}
        \begin{spacing}{1.2}
            \tituloes
        \end{spacing}
        }
        
        \vspace{2cm}
        
        \begin{tabular}{rl}
            \textbf{Autor:} & \alumno \\[0.5cm]
            \textbf{Director:} & \tutor \\
        \end{tabular}
        
        \vfill
        
        {\large Madrid, junio de 2026}
        
        \vspace{1cm}
        
        \rule{\textwidth}{1.5pt}
        
        \vspace{0.5cm}
        
        {\normalsize\facultad}\\[0.2cm]
        {\normalsize\universidad}
        
    \end{center}
\end{titlepage}

%% --- Página en blanco después de portada ---
\newpage
\thispagestyle{empty}
\mbox{}
\newpage

%% --- RESUMEN ---
\chapter*{Resumen}
\addcontentsline{toc}{chapter}{Resumen}

Esto tiene que ser lo último al realizar y son unas 250 - 300 palabras.

\vspace{1cm}

\textbf{Palabras clave:} Inteligencia Artificial, Lesión Renal Aguda, Machine Learning, Predicción Clínica, Datos Hospitalarios, XGBoost.

%% --- ABSTRACT ---
\chapter*{Abstract}
\addcontentsline{toc}{chapter}{Abstract}

\vspace{1cm}

\textbf{Keywords:} Artificial Intelligence, Acute Kidney Injury, Machine Learning, Clinical Prediction, Hospital Data, XGBoost.

%% --- AGRADECIMIENTOS (opcional) ---
\chapter*{Agradecimientos}
\addcontentsline{toc}{chapter}{Agradecimientos}

Sección opcional para expresar agradecimientos a las personas e instituciones que han contribuido al desarrollo de este trabajo.
\clearpage

%% --- ÍNDICE GENERAL ---
\tableofcontents
\clearpage

%% --- ÍNDICE DE FIGURAS ---
\newpage
\listoffigures
\addcontentsline{toc}{chapter}{Índice de Figuras}

%% --- ÍNDICE DE TABLAS ---
\newpage
\listoftables
\addcontentsline{toc}{chapter}{Índice de Tablas}

%% --- ÍNDICE DE CÓDIGO (opcional) ---
%% \lstlistoflistings
%% \addcontentsline{toc}{chapter}{Índice de Código}

%% ==================== CONTENIDO PRINCIPAL ==================== %%
\newpage
\pagenumbering{arabic}

%% --- CAPÍTULO 1: INTRODUCCIÓN ---
\chapter{Introducción}
\label{chap:introduccion}

\section{La lesión renal aguda como problema clínico}

La lesión renal aguda, conocida habitualmente como AKI por sus siglas en inglés (\textit{Acute Kidney Injury}), es una complicación que los profesionales sanitarios encuentran con notable frecuencia en el entorno hospitalario. Se trata de un deterioro brusco de la función renal que altera la capacidad del organismo para eliminar productos de desecho metabólico y mantener el equilibrio de líquidos y electrolitos. Cuando aparece, el manejo del paciente se complica de forma considerable: hay que ajustar las dosis de medicamentos que se eliminan por vía renal, vigilar más estrechamente las analíticas, controlar el balance hídrico y, en los casos más graves, plantear técnicas de depuración extrarrenal como la hemodiálisis.

El riñón es un órgano con funciones que van mucho más allá de la simple producción de orina. Regula la presión arterial, mantiene el equilibrio ácido-base, controla los niveles de potasio y sodio en sangre, y participa en la producción de hormonas como la eritropoyetina. Cuando su función se deteriora de forma aguda, todas estas funciones se ven comprometidas simultáneamente, lo que explica por qué la AKI tiene un impacto tan amplio sobre el estado general del paciente.

Las cifras epidemiológicas dan una idea de la magnitud del problema. Según los estudios más recientes, entre el 10\% y el 15\% de los pacientes hospitalizados desarrollan algún grado de AKI durante su ingreso \parencite{hoste2018epidemiology}. En unidades de cuidados intensivos, donde los pacientes están más graves y expuestos a más factores de riesgo, esta cifra asciende hasta el 50-60\% \parencite{hoste2015akiepi}. El estudio AKI-EPI, que analizó datos de 97 unidades de cuidados intensivos en 33 países, encontró que más de la mitad de los pacientes críticos cumplían criterios de AKI en algún momento de su estancia \parencite{hoste2015akiepi}.

Pero el problema no termina cuando el paciente sale del hospital. La AKI no es simplemente un episodio agudo que se resuelve sin consecuencias. Los pacientes que la desarrollan tienen mayor riesgo de mortalidad tanto a corto como a largo plazo, y existe una conexión bien establecida entre la lesión renal aguda y la progresión hacia enfermedad renal crónica \parencite{chawla2014akickd}. Un riñón que ha sufrido un episodio de daño agudo queda más vulnerable a futuras agresiones, y en algunos casos el deterioro funcional persiste de forma permanente. Esta relación bidireccional entre AKI y enfermedad renal crónica ha llevado a algunos autores a hablar de un síndrome interconectado más que de dos entidades separadas.

Desde una perspectiva de gestión sanitaria, la AKI supone también una carga económica considerable. Los pacientes que la desarrollan requieren estancias más prolongadas, más pruebas diagnósticas, más intervenciones terapéuticas y, en ocasiones, tratamiento sustitutivo renal. Aunque es difícil cuantificar exactamente el coste adicional que supone cada caso de AKI, las estimaciones hablan de incrementos significativos en el gasto hospitalario asociado a esta complicación.

\section{El problema del diagnóstico tardío}

Si la AKI es tan frecuente y tiene consecuencias tan importantes, cabe preguntarse por qué no se detecta antes. La respuesta tiene que ver con las limitaciones de los marcadores diagnósticos disponibles.

El criterio diagnóstico más utilizado para definir la AKI se basa en los cambios en la creatinina sérica. Las guías KDIGO, que constituyen el estándar internacional actual, establecen que existe AKI cuando la creatinina aumenta al menos 0,3 mg/dL en 48 horas, o cuando aumenta al menos 1,5 veces respecto al valor basal en 7 días, o cuando la diuresis desciende por debajo de 0,5 mL/kg/h durante 6 horas \parencite{kdigo2012aki}. De estos criterios, el más utilizado en la práctica clínica y en investigación es el basado en creatinina, porque es el más fácil de medir de forma objetiva.

El problema es que la creatinina es un marcador funcional con un retraso inherente. La creatinina es un producto del metabolismo muscular que se filtra libremente en el glomérulo y se elimina por la orina. Cuando el riñón empieza a fallar, la cantidad de creatinina que se filtra disminuye y, por tanto, su concentración en sangre aumenta. Sin embargo, este aumento no es inmediato. La creatinina tiene que acumularse en el organismo antes de que su elevación sea detectable en una analítica, y este proceso puede tardar horas o incluso días.

Esto significa que cuando el médico ve una creatinina elevada en los análisis del paciente, el daño renal ya lleva tiempo produciéndose. La ventana de oportunidad para actuar de forma preventiva se ha reducido considerablemente. En ese momento, las opciones se limitan a intentar frenar la progresión del daño y evitar que empeore, pero ya no es posible prevenir lo que ya ha ocurrido.

Durante años se han buscado biomarcadores más precoces que pudieran detectar el daño renal antes de que la creatinina empiece a subir. Moléculas como la NGAL (\textit{Neutrophil Gelatinase-Associated Lipocalin}), la KIM-1 (\textit{Kidney Injury Molecule-1}) o la cistatina C han mostrado resultados prometedores en estudios de investigación. Sin embargo, ninguno de estos marcadores ha logrado implantarse de forma generalizada en la práctica clínica habitual, ya sea por cuestiones de coste, disponibilidad o falta de evidencia suficiente sobre su utilidad para guiar decisiones terapéuticas.

Esta situación ha llevado a explorar enfoques alternativos para la detección temprana. Si los biomarcadores individuales tienen limitaciones, quizá sea posible combinar múltiples variables clínicas para identificar patrones de riesgo antes de que el daño se manifieste. Y aquí es donde entran las técnicas de inteligencia artificial.

\section{Inteligencia artificial y predicción clínica}

Los hospitales modernos son, probablemente sin pretenderlo, enormes generadores de datos. Cada ingreso produce un rastro digital que incluye analíticas seriadas, registros de medicación, constantes vitales, diagnósticos codificados, notas clínicas, resultados de pruebas de imagen y mucho más. En un hospital de tamaño medio, hablamos de millones de registros acumulados a lo largo de los años.

La mayor parte de esta información se utiliza exclusivamente para el seguimiento individual del paciente. El médico consulta los análisis de su paciente, revisa la evolución de sus constantes, comprueba qué medicación tiene prescrita. Pero rara vez se explota esta información de forma agregada para detectar patrones que permitan predecir complicaciones antes de que ocurran.

Las técnicas de aprendizaje automático ofrecen precisamente esa posibilidad. Los algoritmos de \textit{machine learning} pueden analizar grandes volúmenes de datos históricos, encontrar relaciones entre variables y construir modelos capaces de estimar el riesgo de un evento futuro a partir de la información disponible en un momento dado. No se trata de sustituir el juicio clínico del médico, sino de proporcionarle una herramienta adicional que le ayude a priorizar su atención hacia los pacientes con mayor riesgo.

En el campo específico de la predicción de AKI, esta aproximación ha generado resultados prometedores. El trabajo más conocido es probablemente el de Tomašev et al., publicado en la revista \textit{Nature} en 2019 \parencite{tomasev2019nature}. Este estudio, desarrollado en colaboración con DeepMind (la división de inteligencia artificial de Google), utilizó datos de más de 700.000 pacientes del sistema de salud del Departamento de Veteranos de Estados Unidos para entrenar un modelo de aprendizaje profundo. Los resultados mostraron que era posible predecir la aparición de AKI hasta 48 horas antes del diagnóstico clínico convencional, con un área bajo la curva ROC de 0,921.

Otros grupos de investigación han obtenido resultados similares utilizando diferentes metodologías y poblaciones. Churpek et al. desarrollaron y validaron un modelo de \textit{machine learning} para predicción de AKI en pacientes hospitalizados, demostrando que el modelo mantenía su rendimiento cuando se validaba en hospitales diferentes a aquellos donde se había entrenado \parencite{churpek2020validation}. Koyner et al. construyeron un modelo específico para pacientes ingresados, utilizando variables disponibles de forma rutinaria en la historia clínica electrónica \parencite{koyner2018ml}. Lin et al. aplicaron Random Forest para predecir la mortalidad en pacientes con AKI en UCI, identificando las variables con mayor capacidad predictiva \parencite{lin2019randomforest}.

Sin embargo, la literatura también es clara en sus advertencias sobre las limitaciones de estos modelos. Las revisiones sistemáticas más recientes han puesto de manifiesto que muchos de los estudios publicados presentan problemas metodológicos importantes \parencite{rank2022mlaki, shi2025aiaki}. Entre los problemas más frecuentes se encuentran la falta de validación externa (el modelo solo se prueba en los mismos datos donde se entrenó), el alto riesgo de sesgo según herramientas como PROBAST, la heterogeneidad en las definiciones de AKI utilizadas y la falta de información sobre la calibración de las probabilidades predichas.

Además, la inmensa mayoría de los estudios publicados proceden de Estados Unidos o del norte de Europa. Las poblaciones de pacientes, los sistemas de salud, los patrones de prescripción farmacológica y las prácticas clínicas pueden ser diferentes en otros contextos. No está claro que un modelo entrenado con datos de hospitales americanos vaya a funcionar igual de bien en un hospital español, con sus propias características organizativas y poblacionales.

Esta situación subraya la necesidad de desarrollar y evaluar modelos en entornos clínicos diversos, incluyendo el contexto sanitario español, que está notablemente infrarrepresentado en la literatura sobre predicción de AKI mediante inteligencia artificial.

\section{Planteamiento del trabajo}

Este Trabajo de Fin de Grado se sitúa en esa línea de investigación aplicada. El objetivo principal es desarrollar y evaluar un modelo de aprendizaje automático capaz de predecir la aparición de lesión renal aguda en las primeras 48 horas de ingreso hospitalario, utilizando datos clínicos reales procedentes de HM Hospitales.

El trabajo se desarrolla en el contexto del Grado en Ingeniería Informática Biomédica de la Universidad Camilo José Cela, y aprovecha la posibilidad de acceder a la base de datos clínica Doctoris utilizada en los hospitales del grupo HM. Esta circunstancia ofrece una oportunidad poco habitual en trabajos académicos de grado: trabajar con datos hospitalarios reales, no simulados ni procedentes de repositorios públicos internacionales.

Trabajar con datos reales introduce complejidades adicionales frente a los \textit{datasets} de laboratorio. La información clínica presenta valores ausentes, codificaciones heterogéneas, registros incompletos y particularidades derivadas de la práctica asistencial cotidiana. No todos los pacientes tienen todas las analíticas, no todas las prescripciones están correctamente codificadas, no todos los campos están cumplimentados de la misma manera. Esta realidad obliga a dedicar una parte sustancial del trabajo a la fase de ingeniería del dato: extracción, limpieza, transformación y validación de la información antes de poder utilizarla para entrenar modelos.

El enfoque metodológico adoptado prioriza la interpretabilidad sobre la complejidad. Aunque existen arquitecturas más sofisticadas, como las redes neuronales profundas utilizadas por Tomašev et al. \parencite{tomasev2019nature}, estas requieren volúmenes de datos considerablemente mayores y plantean retos de interpretabilidad que dificultan su aceptación en entornos clínicos. Para una prueba de concepto desarrollada en el contexto de un TFG, con las limitaciones de tamaño muestral inherentes a una cohorte monocéntrica, resulta más adecuado emplear modelos cuyo funcionamiento interno pueda analizarse y cuyas predicciones puedan explicarse en términos clínicamente comprensibles.

Por este motivo, los modelos seleccionados son la regresión logística y el Random Forest. La regresión logística es un modelo lineal ampliamente utilizado en investigación biomédica, cuya principal ventaja es la facilidad para interpretar los coeficientes en términos de \textit{odds ratios}. El Random Forest es un modelo de ensamble basado en árboles de decisión, capaz de capturar relaciones no lineales entre variables sin necesidad de especificarlas explícitamente. La comparación entre ambos permite contrastar el rendimiento de un enfoque interpretable y paramétrico frente a otro más flexible y no lineal.

La variable objetivo se define como el desarrollo de AKI en las 48 horas siguientes al valor basal de creatinina del ingreso, adaptando el criterio de incremento absoluto de creatinina de las guías KDIGO \parencite{kdigo2012aki}. Como predictores se utilizan variables disponibles de forma rutinaria en el momento del ingreso: variables demográficas (edad, sexo), variables analíticas (creatinina basal, sodio, potasio) y variables farmacológicas (exposición a AINEs, IECAs, ARA-II, diuréticos, vancomicina y aminoglucósidos).

Es importante señalar desde el principio que este trabajo constituye una prueba de concepto académica, no una herramienta lista para implantación clínica inmediata. El objetivo no es desarrollar un sistema de alertas que pueda ponerse en producción mañana en un hospital, sino demostrar la viabilidad técnica del enfoque, evaluar el rendimiento de modelos \textit{baseline} sobre datos reales, identificar qué variables parecen más relevantes para la predicción y sentar las bases para posibles desarrollos futuros más ambiciosos.

\section{Estructura del documento}

El presente documento se organiza en nueve capítulos, además de las referencias bibliográficas y los anexos. A continuación se describe brevemente el contenido de cada uno de ellos.

El \textbf{Capítulo 2} presenta el marco teórico del trabajo. Se revisan en profundidad la definición y los criterios diagnósticos de la AKI según las guías KDIGO, la epidemiología de esta complicación tanto a nivel global como en el contexto hospitalario, los principales factores de riesgo con especial atención a la nefrotoxicidad farmacológica, y el estado del arte en predicción de AKI mediante técnicas de inteligencia artificial. También se presentan los fundamentos de los modelos de clasificación utilizados (regresión logística y Random Forest), las consideraciones especiales sobre evaluación de modelos en problemas con clases desbalanceadas, y la importancia de la interpretabilidad en modelos destinados a aplicaciones clínicas.

El \textbf{Capítulo 3} expone la justificación del trabajo. Se argumenta la relevancia clínica del problema abordado, la oportunidad tecnológica que representan los datos hospitalarios estructurados y las técnicas de aprendizaje automático, y el contexto institucional en que se desarrolla el proyecto.

El \textbf{Capítulo 4} detalla los objetivos del trabajo, tanto el objetivo general como los objetivos específicos que guían el desarrollo del proyecto.

El \textbf{Capítulo 5} describe la metodología empleada. Se presenta el entorno técnico utilizado (SQL Server, Python, Jupyter Notebook, librerías científicas), se explica el proceso de definición de la cohorte de estudio, se detalla la construcción de la variable objetivo, se describen las variables predictoras incluidas en el modelo, se expone el preprocesamiento de datos realizado y se presenta la estrategia de modelado adoptada.

El \textbf{Capítulo 6} presenta el desarrollo del trabajo propiamente dicho. Se describe el pipeline de análisis implementado, el análisis exploratorio inicial de los datos, las decisiones de ingeniería de variables adoptadas y el proceso de entrenamiento de los modelos.

El \textbf{Capítulo 7} recoge los resultados obtenidos y su discusión. Se presentan las métricas de rendimiento de los modelos entrenados, el análisis del umbral de decisión, la interpretación de las variables más relevantes, una demostración del funcionamiento del modelo sobre casos simulados, y una discusión que incluye la interpretación clínica de los resultados, la comparación con la literatura existente y las limitaciones del estudio.

El \textbf{Capítulo 8} presenta las conclusiones principales derivadas del trabajo realizado, sintetizando los hallazgos más relevantes y valorando el grado de cumplimiento de los objetivos planteados.

El \textbf{Capítulo 9} propone líneas de trabajo futuro que podrían dar continuidad a este proyecto, identificando áreas de mejora y posibles extensiones.

Finalmente, se incluyen las \textbf{referencias bibliográficas} citadas a lo largo del documento y varios \textbf{anexos} con información complementaria: un diccionario de variables con la descripción detallada de cada campo utilizado, el código SQL empleado para la extracción de la cohorte, y fragmentos de código Python correspondientes a las fases más relevantes del análisis.

%% --- CAPÍTULO 2: MARCO TEÓRICO ---
\chapter{Marco Teórico}
\section{Lesión renal aguda: definición y criterios diagnósticos}

La lesión renal aguda, conocida por sus siglas en inglés AKI (Acute Kidney Injury), es un síndrome clínico que se caracteriza por una disminución brusca de la función renal. Durante décadas se empleó el término "insuficiencia renal aguda", pero la comunidad nefrológica ha ido abandonándolo progresivamente porque transmite la idea de un fallo completo del órgano, cuando en realidad muchos pacientes presentan deterioros más sutiles que también tienen consecuencias clínicas relevantes \parencite{kdigo2012aki}.

Uno de los problemas históricos en este campo ha sido la falta de una definición estandarizada. Hasta principios de los años 2000, la literatura recogía más de 30 definiciones distintas de lo que constituía un "fallo renal agudo", lo que dificultaba enormemente comparar estudios entre sí. Esta situación cambió en 2004, cuando el grupo ADQI (Acute Dialysis Quality Initiative) propuso los criterios RIFLE, un acrónimo que recoge cinco categorías de gravedad creciente: Riesgo, Injuria, Fallo, Pérdida y Enfermedad renal terminal. Por primera vez se disponía de un sistema de clasificación basado en cambios objetivos en la creatinina sérica y en el volumen de orina.

Tres años después, la red AKIN (Acute Kidney Injury Network) introdujo una modificación importante: añadió como criterio diagnóstico un incremento absoluto de creatinina de al menos 0,3 mg/dL en un periodo de 48 horas. Este umbral, aparentemente pequeño, se incluyó tras comprobarse en varios estudios que elevaciones tan modestas ya se asociaban con peores resultados clínicos, incluyendo mayor mortalidad hospitalaria \parencite{khwaja2012kdigo}(Khwaja, 2012).

La síntesis definitiva llegó en 2012 con las guías KDIGO (Kidney Disease: Improving Global Outcomes), que son las que se utilizan en la actualidad y las que he empleado en este trabajo. Según estos criterios, se diagnostica AKI cuando se cumple cualquiera de las siguientes condiciones:

- Un aumento de la creatinina sérica igual o superior a 0,3 mg/dL en un plazo de 48 horas.
- Un aumento de la creatinina sérica igual o superior a 1,5 veces el valor basal, presumiblemente ocurrido en los 7 días previos.
- Un volumen de orina inferior a 0,5 mL/kg/hora durante al menos 6 horas.

Las guías KDIGO también establecen tres estadios de gravedad. El estadio 1 corresponde al cumplimiento mínimo de los criterios; el estadio 2 implica un aumento de creatinina entre 2 y 3 veces el valor basal; y el estadio 3 se reserva para incrementos superiores a 3 veces el basal, creatininas por encima de 4 mg/dL con un aumento agudo de al menos 0,5 mg/dL, o la necesidad de terapia de reemplazo renal.

Un aspecto que merece atención es la determinación del valor basal de creatinina. En la práctica clínica, no siempre se dispone de una analítica previa del paciente, lo que obliga a recurrir a estimaciones. Las guías sugieren que, en ausencia de datos, se puede asumir una tasa de filtrado glomerular de 75 mL/min/1,73 m² y calcular hacia atrás la creatinina esperada según edad y sexo. Sin embargo, este método tiene limitaciones evidentes y puede infraestimar la incidencia real de AKI en poblaciones con enfermedad renal crónica subyacente no diagnosticada \parencite{palevsky2013kdoqi}.

En el presente trabajo, he optado por definir la AKI utilizando el criterio de incremento absoluto de creatinina. Concretamente, se considera que un paciente desarrolla AKI en las primeras 48 horas de ingreso si su creatinina máxima en ese periodo supera en 0,3 mg/dL o más a la creatinina basal registrada al ingreso. Esta elección se debe a que los datos disponibles en la base Doctoris no incluyen información fiable sobre diuresis horaria, y el criterio de 0,3 mg/dL es el que mejor se adapta a una ventana temporal corta de predicción.

\section{Epidemiología y relevancia clínica de la AKI hospitalaria}
La AKI no es una rareza clínica. Más bien al contrario: se trata de una de las complicaciones más frecuentes en el entorno hospitalario, y su incidencia ha ido aumentando en las últimas décadas. Este incremento se explica, al menos en parte, por el envejecimiento de la población, el aumento de comorbilidades como la diabetes o la insuficiencia cardíaca, y el uso más extendido de procedimientos diagnósticos y terapéuticos con potencial nefrotóxico.

Según la revisión de Hoste y colaboradores publicada en Nature Reviews Nephrology, las estimaciones de prevalencia de AKI varían enormemente dependiendo del contexto clínico y los criterios utilizados. En hospitalizaciones generales, las cifras oscilan entre el 10 y el 15 por ciento de los ingresos. En unidades de cuidados intensivos, la situación es bastante más preocupante: estudios multicéntricos como el AKI-EPI, que incluyó cerca de 2.000 pacientes de 97 centros en distintos países, encontraron que más de la mitad de los pacientes críticos desarrollaban algún grado de AKI durante la primera semana de ingreso en UCI \parencite{hoste2015akiepi}. Son cifras que invitan a la reflexión.

La mortalidad asociada a la AKI es considerable. En pacientes hospitalizados, el desarrollo de AKI multiplica el riesgo de muerte intrahospitalaria, con tasas que pueden superar el 20-25 por ciento en formas moderadas y alcanzar el 50 por ciento o más cuando se requiere diálisis. Un dato que resulta muy llamativo es que incluso los estadios más leves de AKI, aquellos en los que la creatinina apenas sube unas décimas, ya se asocian con un aumento significativo de la mortalidad. Esto refuerza la idea de que no estamos ante un problema menor o transitorio.

Más allá de la fase aguda, la AKI tiene consecuencias a largo plazo que durante mucho tiempo pasaron relativamente desapercibidas. Como señalan Chawla y colaboradores en su revisión del \textbf{New England Journal of Medicine}, la AKI y la enfermedad renal crónica no son entidades separadas sino que están íntimamente conectadas. Un episodio de AKI, aunque se resuelva aparentemente por completo, deja huella: aumenta el riesgo de desarrollar enfermedad renal crónica en el futuro, acelera la progresión de la ERC preexistente, y eleva la mortalidad cardiovascular a medio y largo plazo \parencite{chawla2014akickd}.

El impacto económico tampoco es despreciable. Los pacientes que desarrollan AKI permanecen más tiempo hospitalizados, requieren más recursos y generan costes significativamente mayores. En algunos sistemas sanitarios se ha estimado que la AKI puede incrementar el coste de una hospitalización en varios miles de euros. A esto hay que añadir los costes derivados del seguimiento posterior, las posibles necesidades de diálisis crónica y la pérdida de productividad.

Todo lo anterior subraya una idea central: la detección precoz de la AKI tiene un valor clínico indiscutible. Si se pudiera identificar a los pacientes con mayor riesgo de desarrollar daño renal antes de que este se manifieste analíticamente, existiría una ventana de oportunidad para intervenir. Las intervenciones preventivas no son complicadas: optimizar la hidratación, evitar o ajustar fármacos nefrotóxicos, monitorizar más estrechamente la función renal. El problema es saber a quién aplicar estas medidas con antelación suficiente.

En el contexto de HM Hospitales, donde se ha desarrollado este trabajo, no existían hasta la fecha datos específicos sobre la incidencia de AKI. El análisis de la cohorte extraída de la base de datos Doctoris reveló una prevalencia de AKI en las primeras 48 horas del 2,5 por ciento aproximadamente. Esta cifra es inferior a las reportadas en series de UCI, lo cual tiene sentido dado que la población estudiada incluye pacientes de hospitalización general. De todos modos, incluso este porcentaje aparentemente bajo representa un número absoluto considerable de pacientes si se extrapola al volumen total de ingresos anuales de la red hospitalaria.
\section{Factores de riesgo de AKI y el papel de la nefrotoxicidad farmacológica}

El desarrollo de una lesión renal aguda rara vez obedece a una causa única. Lo habitual es que confluyan varios factores que, actuando de forma conjunta, sobrepasan la capacidad de adaptación del riñón. Conocer estos factores de riesgo resulta fundamental tanto para la prevención como para el desarrollo de modelos predictivos.

Entre los factores no modificables, la edad ocupa un lugar destacado. El riñón envejecido tiene menor reserva funcional: pierde nefronas, presenta mayor fibrosis intersticial y responde peor a las agresiones. Por eso no sorprende que la incidencia de AKI aumente de forma marcada a partir de los 65 años. En la cohorte analizada en este trabajo, la variable "Edad mayor o igual a 65" se incluyó precisamente para capturar este efecto.

La presencia de enfermedad renal crónica previa es otro factor de riesgo bien establecido. Un paciente que ya parte de una función renal disminuida tiene menos margen antes de cumplir criterios de AKI, y además su riñón es más vulnerable a nuevas agresiones. Algo similar ocurre con la diabetes mellitus, la insuficiencia cardíaca y la cirrosis hepática: todas ellas predisponen al desarrollo de daño renal agudo por mecanismos diversos.

Sin embargo, los factores que más interés tienen desde el punto de vista preventivo son los modificables. Y entre ellos, la exposición a fármacos nefrotóxicos merece una atención especial. Se estima que los medicamentos están implicados en aproximadamente el 20-30 por ciento de los casos de AKI hospitalaria. En algunos contextos, como las unidades de cuidados intensivos, esta proporción puede ser aún mayor.

Los antiinflamatorios no esteroideos, los famosos AINEs, son probablemente los nefrotóxicos más extendidos. Su mecanismo de daño renal está bien caracterizado: al inhibir la ciclooxigenasa, bloquean la síntesis de prostaglandinas renales que normalmente mantienen la vasodilatación de la arteriola aferente. En condiciones de perfusión renal comprometida (deshidratación, insuficiencia cardíaca, uso concomitante de diuréticos), esta inhibición puede precipitar una caída brusca del filtrado glomerular. El riesgo es especialmente alto cuando se combinan AINEs con inhibidores del sistema renina-angiotensina y diuréticos, una asociación conocida en la literatura anglosajona como "triple whammy" \parencite{zhang2017nsaid}.

Los aminoglucósidos constituyen otro grupo clásico de nefrotóxicos. Estos antibióticos, que incluyen gentamicina, amikacina y tobramicina, se acumulan en las células del túbulo proximal a través de un mecanismo de endocitosis mediado por el receptor megalina. Una vez dentro de la célula, alteran la función mitocondrial y lisosomal, lo que eventualmente conduce a la muerte celular. La nefrotoxicidad por aminoglucósidos es dosis-dependiente y se relaciona también con la duración del tratamiento. Por este motivo, las pautas de administración en dosis única diaria han ganado terreno frente a las pautas fraccionadas tradicionales \parencite{perazella2022drugaki}.

La vancomicina, un glucopéptido ampliamente utilizado frente a infecciones por grampositivos resistentes, también se asocia con nefrotoxicidad. Durante años existió debate sobre si el daño renal atribuido a vancomicina no se debía en realidad a factores confundentes, pero estudios más recientes han confirmado que el fármaco puede causar daño tubular directo, especialmente cuando se alcanzan concentraciones plasmáticas elevadas o se prolonga el tratamiento \parencite{pais2020vancomycin}. La monitorización de niveles plasmáticos es práctica habitual para minimizar este riesgo.

Los inhibidores de la enzima convertidora de angiotensina (IECAs) y los antagonistas de los receptores de angiotensina II (ARA-II) ocupan un lugar peculiar en esta lista. No son nefrotóxicos en el sentido clásico del término; de hecho, a largo plazo ejercen un efecto nefroprotector en pacientes con proteinuria. Sin embargo, su mecanismo de acción, que consiste en dilatar la arteriola eferente glomerular, puede precipitar una caída del filtrado glomerular en situaciones de hipoperfusión renal. Por eso se recomienda suspenderlos temporalmente cuando un paciente presenta deshidratación, hipotensión o enfermedad aguda intercurrente.

Los diuréticos, especialmente los de asa como la furosemida, contribuyen al riesgo de AKI principalmente a través de la depleción de volumen. Un paciente excesivamente diurizado pierde volumen intravascular efectivo, lo que compromete la perfusión renal. El efecto se magnifica cuando se combinan con AINEs o inhibidores del sistema renina-angiotensina.

En el modelo desarrollado en este trabajo, se incluyeron variables binarias para cada uno de estos grupos farmacológicos: AINE, IECA, ARA-II, diurético, vancomicina y aminoglucósido. Además, se creó una variable derivada que cuenta el número total de fármacos de riesgo que recibe cada paciente, partiendo de la hipótesis de que la exposición simultánea a múltiples nefrotóxicos incrementa el riesgo de forma aditiva o incluso sinérgica.

\section{Predicción de AKI mediante inteligencia artificial: estado del arte}

La idea de anticiparse al desarrollo de la AKI no es nueva. Durante décadas se han propuesto diversos scores clínicos y modelos de riesgo, algunos específicos para contextos concretos como la cirugía cardíaca o la exposición a contraste yodado. Sin embargo, estos modelos tradicionales, basados generalmente en regresión logística con un número reducido de variables, han mostrado un rendimiento modesto en la práctica clínica habitual.

El auge de las técnicas de aprendizaje automático ha abierto nuevas posibilidades en este campo. Los algoritmos de machine learning pueden manejar un número elevado de variables, detectar interacciones no lineales y adaptarse a patrones complejos que difícilmente capturaría un modelo estadístico convencional. No se extraña, por tanto, que en la última década haya proliferado la investigación sobre predicción de AKI mediante inteligencia artificial.

Varias revisiones sistemáticas han intentado sintetizar esta literatura creciente. El trabajo de Rank y colaboradores, publicado en \parencite{rank2022mlaki}, identificó 46 estudios que desarrollaban modelos de machine learning para predecir AKI. Los algoritmos más utilizados fueron Random Forest, gradient boosting y redes neuronales, aunque la regresión logística seguía apareciendo con frecuencia como modelo de referencia para comparación. En cuanto a las variables predictoras, las más recurrentes fueron la edad, la creatinina sérica, el nitrógeno ureico en sangre, el potasio y diversos indicadores de comorbilidad. 

Una revisión más reciente, la de Zhao y colaboradores publicada en \parencite{zhao2024akiml}, analizó específicamente los modelos de ML aplicados al pronóstico de pacientes con AKI ya establecida. Los resultados del meta-análisis mostraron que los métodos de machine learning superaban a los enfoques tradicionales en la predicción de mortalidad, con un área bajo la curva ROC de 0,831 frente a 0,772 para los modelos clásicos. Esta diferencia, aunque estadísticamente significativa, da a una reflexión: en los modelos de machine learning existe cierta mejora , aunque la diferncia no es tan grande comopodría esperarse leyendo ciertos estudios.

El estudio que probablemente más repercusión mediática ha tenido en este campo fue el publicado por Tomašev y colaboradores en la revista \parencite{tomasev2019nature}. Utilizando datos de más de 700.000 pacientes del sistema de salud de veteranos de Estados Unidos, el equipo de DeepMind (agencia de Google dedicada a inteligencia artificial) desarrolló un modelo de redes neuronales recurrentes capaz de predecir AKI con hasta 48 horas de antelación, alcanzando un área bajo la curva de 0,921. El trabajo generó considerable expectación, aunque también críticas metodológicas: el modelo funcionaba mucho mejor en hombres que en mujeres (la cohorte de veteranos era predominantemente masculina), y las tasas de falsos positivos eran elevadas, lo que planteaba dudas sobre su utilidad práctica.

Dejando a un lado los titulares más optimistas, lo que muestran las revisiones es un escenario con cosas positivas pero también con bastantes limitaciones. Por el lado positivo, existe evidencia de que los modelos de ML pueden alcanzar rendimientos aceptables e incluso buenos en la predicción de AKI. Por el lado menos favorable, la mayoría de los estudios presentan limitaciones metodológicas importantes. La revisión de Shi y colaboradores, publicada en \parencite{shi2025aiaki}, evaluó 47 estudios utilizando las herramientas PROBAST y TRIPOD y encontró que 44 de ellos tenían alto riesgo de sesgo. Los problemas más frecuentes eran la falta de validación externa (solo 14 estudios la realizaban), el uso de cohortes muy seleccionadas, y deficiencias en el reporte de la metodología.

Esta última observación tiene implicaciones directas para el presente trabajo. Desarrollar un modelo predictivo que funcione bien en los datos con los que se entrena es relativamente sencillo; lo difícil es que ese modelo generalice a poblaciones diferentes. Por eso, cualquier trabajo de estas características debe ser transparente respecto a sus limitaciones y evitar conclusiones excesivamente optimistas.

Otra cuestión relevante es la brecha entre el desarrollo de modelos y su implementación clínica real. La mayoría de los estudios publicados se quedan en la fase de desarrollo y validación interna. Son muy pocos los que llegan a integrarse en sistemas de información hospitalaria o a evaluarse en ensayos prospectivos que midan su impacto en resultados de salud. Esta brecha no es exclusiva del campo de la AKI, pero resulta particularmente llamativa dado el volumen de publicaciones existentes.

En definitiva, el estado del arte sugiere que la predicción de AKI mediante machine learning es una línea de investigación prometedora pero todavía inmadura. Los modelos publicados muestran capacidad discriminativa razonable, pero su traslación a la práctica clínica requiere abordar cuestiones de validación externa, interpretabilidad, integración en flujos de trabajo y, en último término, demostración de beneficio para los pacientes.

\section{Modelos de aprendizaje automático para clasificación clínica}

En este trabajo se han empleado dos algoritmos de clasificación: la regresión logística y el Random Forest. La elección de estos dos modelos no es casual. Representan, en cierto modo, dos filosofías distintas de abordar los problemas de predicción, y su comparación permite obtener perspectivas complementarias.

\subsection{Regresión logística}

La regresión logística es probablemente el modelo de clasificación más utilizado en investigación clínica y epidemiológica. Su popularidad se debe a una combinación de factores: es matemáticamente sencilla, sus resultados son directamente interpretables, y cuenta con décadas de uso que la han convertido en un estándar aceptado por clínicos y revisores de revistas médicas.

El modelo parte de una combinación lineal de las variables predictoras, que luego se transforma mediante la función logística para obtener una probabilidad entre 0 y 1. En términos formales, si llamamos "p" a la probabilidad de que ocurra el evento de interés (en nuestro caso, desarrollar AKI), el modelo estima:

\begin{equation}
\log\left(\frac{p}{1-p}\right) = \beta_0 + \beta_1 X_1 + \beta_2 X_2 + \cdots + \beta_n X_n
\end{equation}

Los coeficientes $\beta$ representan el cambio en el logaritmo de odds por cada unidad de incremento en la variable correspondiente. Exponenciando estos coeficientes se obtienen los odds ratios, que tienen una interpretación intuitiva: un OR de 2 para una variable indica que la presencia de esa característica duplica las odds de desarrollar el evento.

Esta interpretabilidad constituye la principal fortaleza de la regresión logística en el contexto clínico. Un médico puede entender que un paciente con creatinina basal elevada tiene mayor riesgo de AKI, y puede cuantificar aproximadamente cuánto mayor es ese riesgo. Esto facilita la aceptación del modelo y su uso en la toma de decisiones.

Sin embargo, la regresión logística también tiene limitaciones. Asume una relación lineal entre los predictores y el logaritmo de odds, lo cual no siempre se cumple en la realidad. Por ejemplo, la relación entre edad y riesgo de AKI podría ser no lineal, con un incremento más pronunciado a partir de cierta edad. Además, el modelo no captura de forma automática las interacciones entre variables: si el efecto de los AINEs es particularmente nocivo en pacientes que también toman diuréticos, la regresión logística básica no detectará esta interacción a menos que se especifique explícitamente \parencite{hosmer2013logistic}.

\subsection{Random Forest}

El Random Forest, propuesto por Leo Breiman en 2001 \parencite{breiman2001rf}, representa un enfoque radicalmente distinto. En lugar de ajustar una única ecuación, el algoritmo construye un conjunto (o "bosque") de árboles de decisión, cada uno entrenado con una muestra aleatoria de los datos y utilizando un subconjunto aleatorio de variables en cada división. La predicción final se obtiene agregando las predicciones individuales de todos los árboles, típicamente mediante votación mayoritaria en problemas de clasificación.

Esta estrategia de ensemble tiene varias ventajas. En primer lugar, reduce el riesgo de sobreajuste que aqueja a los árboles de decisión individuales. Un árbol solo, si se deja crecer sin restricciones, puede aprender de memoria las peculiaridades de los datos de entrenamiento sin capturar patrones generalizables. Al promediar muchos árboles construidos de forma ligeramente diferente, el Random Forest suaviza este problema.

En segundo lugar, el algoritmo es capaz de modelar relaciones no lineales e interacciones de forma automática, sin necesidad de especificarlas previamente. Cada árbol puede dividir el espacio de predictores de formas complejas, y el conjunto captura patrones que un modelo lineal pasaría por alto.

En tercer lugar, el Random Forest proporciona medidas de importancia de variables. Existen varias formas de calcularlas, siendo las más comunes la importancia por permutación (cuánto empeora el rendimiento del modelo si se barajan aleatoriamente los valores de una variable) y la importancia basada en la disminución de impureza de Gini. Estas medidas permiten identificar qué predictores son más relevantes para las predicciones del modelo.

El principal inconveniente del Random Forest es su menor interpretabilidad. No produce coeficientes ni odds ratios; la relación entre cada variable y el resultado está distribuida a lo largo de cientos de árboles y no puede resumirse en una fórmula sencilla. Esto dificulta explicar a un clínico exactamente por qué el modelo ha asignado determinada probabilidad a un paciente concreto.

\subsection{¿Cuál es mejor?} 

La comparación sistemática entre Random Forest y regresión logística ha sido objeto de numerosos estudios. El trabajo de Couronné y colaboradores, publicado en *BMC Bioinformatics* en 2018 \parencite{couronne2018rflr}, analizó el rendimiento de ambos algoritmos en 243 datasets reales de características variadas. Sus resultados mostraron que el Random Forest superaba a la regresión logística en aproximadamente el 69\% de los casos, con una diferencia media en AUC de 0,04 puntos. La ventaja del Random Forest tendía a ser mayor en datasets con muchas variables y relaciones no lineales.

Sin embargo, esta superioridad no es universal. En problemas con pocas variables, muestras pequeñas o relaciones aproximadamente lineales, la regresión logística puede igualar o incluso superar al Random Forest, con la ventaja añadida de su mayor transparencia. Por eso tiene sentido, como se hace en este trabajo, entrenar y comparar ambos modelos en lugar de asumir de entrada cuál funcionará mejor.

\section{Evaluación de modelos en problemas con clases desbalanceadas}

Un aspecto metodológico que merece atención especial en este trabajo es el desequilibrio entre clases. De los 1.175 ingresos hospitalarios analizados, solo 29 desarrollaron AKI en las primeras 48 horas, lo que representa aproximadamente un 2,5\% del total. Este tipo de desbalanceo es habitual en problemas de predicción clínica: los eventos adversos que queremos predecir suelen ser, por definición, relativamente infrecuentes.

El problema con los datos desbalanceados es que las métricas de evaluación convencionales pueden resultar engañosas. Consideremos la exactitud (accuracy), que mide el porcentaje de predicciones correctas. Un modelo que simplemente clasificara a todos los pacientes como "no AKI" obtendría una exactitud del 97,5\%: acertaría en 1.146 de 1.175 casos. Sin embargo, este modelo sería completamente inútil para el propósito clínico de identificar pacientes en riesgo, ya que no detectaría ni un solo caso positivo.

Por este motivo, en problemas desbalanceados conviene recurrir a métricas que evalúen por separado el rendimiento en cada clase. La sensibilidad (también llamada recall o tasa de verdaderos positivos) mide qué proporción de los casos realmente positivos el modelo consigue identificar. En nuestro contexto, una sensibilidad del 70\% significaría que el modelo detecta correctamente 7 de cada 10 pacientes que desarrollan AKI. Esta métrica es particularmente relevante cuando el coste de pasar por alto un caso positivo es alto, como ocurre en muchas aplicaciones médicas.

La precisión (o valor predictivo positivo) adopta la perspectiva complementaria: de todos los pacientes que el modelo clasifica como positivos, qué proporción realmente lo son. Una precisión baja indica que el modelo genera muchas alarmas falsas, lo cual también tiene consecuencias prácticas: si un sistema de alerta hospitalario produce demasiados falsos positivos, el personal sanitario acabará ignorándolo.

El F1-score combina sensibilidad y precisión en una única métrica mediante su media armónica. Resulta útil cuando se quiere un balance entre ambas, aunque tiene la limitación de ponderar ambas por igual, lo que puede no reflejar las prioridades clínicas reales. En algunas situaciones tiene más sentido usar variantes como el F2-score, que da mayor peso a la sensibilidad.

El área bajo la curva ROC (AUC-ROC) ofrece una perspectiva diferente. La curva ROC representa la sensibilidad frente a la tasa de falsos positivos para todos los posibles umbrales de decisión. El AUC resume esta información en un único número entre 0 y 1, donde 0,5 corresponde a un clasificador aleatorio y 1 a un clasificador perfecto. El AUC tiene la ventaja de ser independiente del umbral elegido y de ser relativamente robusto al desbalanceo de clases.

No obstante, cuando el desbalanceo es muy pronunciado, el AUC-ROC puede ofrecer una visión excesivamente optimista. Esto se debe a que la tasa de falsos positivos (eje X de la curva) está dominada por la clase mayoritaria, que es muy numerosa. Algunos autores recomiendan en estos casos prestar más atención a la curva precisión-recall y su área bajo la curva (PR-AUC), que se centra exclusivamente en el rendimiento sobre la clase minoritaria \parencite{saito2015prroc}.

Un aspecto que a menudo se pasa por alto es la elección del umbral de decisión. Por defecto, muchos algoritmos asignan la clase positiva cuando la probabilidad predicha supera 0,5. Sin embargo, este umbral no es óptimo en problemas desbalanceados. Reducirlo (por ejemplo, a 0,2 o 0,3) aumenta la sensibilidad a costa de la precisión, lo cual puede ser apropiado si el objetivo prioritario es no dejar escapar casos positivos. En este trabajo se ha explorado el efecto de variar el umbral para entender mejor el comportamiento de los modelos.

Por último, existen estrategias para abordar el desbalanceo durante el entrenamiento del modelo. El remuestreo de la clase minoritaria (oversampling), el submuestreo de la clase mayoritaria (undersampling), y técnicas sintéticas como SMOTE son opciones comunes. Otra alternativa es asignar pesos diferentes a las clases en la función de pérdida, penalizando más los errores sobre la clase minoritaria. En este trabajo se ha optado por esta última estrategia, utilizando el parámetro class\_weight='balanced' disponible en las implementaciones de scikit-learn.


\section{Interpretabilidad en modelos clínicos}

La interpretabilidad de los modelos predictivos es un tema que ha ganado protagonismo en los últimos años, especialmente en el ámbito sanitario. No basta con que un modelo tenga buen rendimiento; también es deseable, y en muchos contextos imprescindible, entender cómo llega a sus predicciones.

Esta exigencia de transparencia tiene varias justificaciones. En primer lugar, los profesionales sanitarios necesitan confiar en las herramientas que utilizan. Un modelo opaco, tipo "caja negra", que emite predicciones sin explicación alguna, genera resistencia comprensible. El clínico puede preguntarse: ¿por qué debería creer lo que dice este algoritmo? Si no puedo entender su razonamiento, ¿cómo puedo valorar si tiene sentido en este paciente concreto?

En segundo lugar, la interpretabilidad facilita la detección de errores. Un modelo puede tener buen rendimiento aparente pero basarse en características espurias o en artefactos de los datos. Por ejemplo, se han documentado casos de modelos radiológicos que aprendían a distinguir imágenes según el equipamiento utilizado en lugar de las características clínicas relevantes. Sin interpretabilidad, estos problemas pueden pasar desapercibidos.

En tercer lugar, entender qué variables son importantes para el modelo proporciona información clínicamente valiosa. Si un modelo de predicción de AKI da mucho peso a la creatinina basal y a la exposición a AINEs, esto refuerza el conocimiento existente y sugiere dónde enfocar las intervenciones preventivas.

Los modelos de regresión logística tienen interpretabilidad intrínseca. Los coeficientes indican directamente la magnitud y dirección del efecto de cada variable. Un coeficiente positivo significa que valores más altos de esa variable se asocian con mayor probabilidad del evento; el tamaño del coeficiente (o, más intuitivamente, el odds ratio) cuantifica la fuerza de esa asociación.

El Random Forest, en cambio, requiere técnicas de interpretabilidad post-hoc. La más utilizada es la importancia de variables, que puede calcularse de varias formas. La importancia por permutación mide cuánto empeora el rendimiento del modelo cuando se destruye aleatoriamente la información de una variable; variables más importantes producen mayor deterioro. La importancia basada en Gini mide la contribución promedio de cada variable a la reducción de impureza en los nodos de los árboles.

Estas medidas de importancia proporcionan un ranking de las variables más influyentes, pero no indican la dirección del efecto ni permiten cuantificarlo de forma tan directa como los coeficientes de regresión. Para subsanar esto, se han desarrollado métodos como los gráficos de dependencia parcial, que muestran cómo varía la predicción esperada a medida que cambia una variable manteniendo las demás constantes. También existen técnicas más sofisticadas como SHAP (SHapley Additive exPlanations), que descomponen la predicción de cada caso individual en contribuciones atribuibles a cada variable.

En el presente trabajo, la interpretación de los modelos se ha abordado examinando los coeficientes de la regresión logística y la importancia de variables del Random Forest. Esto permite identificar qué predictores son más relevantes según cada modelo y comparar si ambos coinciden en señalar las mismas variables como importantes. Como se verá en la sección de resultados, existe cierta concordancia pero también diferencias interesantes que merecen discusión.

%% --- CAPÍTULO 3: JUSTIFICACIÓN ---
\chapter{Justificación}
\label{chap:justificacion}
El desarrollo de este Trabajo de Fin de Grado responde a una confluencia de factores clínicos, tecnológicos y académicos que hacen del problema abordado un campo de estudio con alto interés tanto para la investigación biomédica como para la práctica asistencial.

\section{Relevancia clínica del problema}
\label{sec:relevancia}
La lesión renal aguda representa una de las complicaciones más frecuentes y con mayor impacto pronóstico dentro del entorno hospitalario. Los datos epidemiológicos disponibles indican que entre el 10\% y el 15\% de los pacientes hospitalizados desarrollan algún grado de AKI durante su ingreso, cifra que asciende hasta el 50-60\% en unidades de cuidados intensivos \parencite{hoste2018epidemiology, hoste2015akiepi}. Las consecuencias de esta complicación trascienden el episodio agudo: los pacientes que desarrollan AKI presentan mayor mortalidad intrahospitalaria, estancias más prolongadas y un riesgo incrementado de progresar hacia enfermedad renal crónica a medio y largo plazo \parencite{chawla2014akickd}.

Desde una perspectiva de gestión sanitaria, la AKI supone también una carga económica considerable. La necesidad de tratamiento sustitutivo renal, la prolongación de los ingresos y las complicaciones asociadas generan costes que podrían reducirse si se actuara de forma preventiva sobre los pacientes con mayor riesgo. Este aspecto adquiere especial relevancia en un contexto de recursos sanitarios limitados, donde la optimización de la atención resulta prioritaria.

El problema central radica en que la detección de la AKI mediante marcadores convencionales, fundamentalmente la creatinina sérica, suele producirse cuando el daño renal ya está establecido. La creatinina es un marcador funcional con un retraso inherente: sus niveles no aumentan de forma inmediata tras la agresión renal, sino cuando la capacidad de filtración glomerular ya se ha deteriorado significativamente. Esta latencia reduce el margen de maniobra para implementar medidas preventivas o correctoras, como el ajuste de fármacos nefrotóxicos, la optimización hemodinámica o la intensificación de la monitorización analítica.

\section{Oportunidad tecnológica}
\label{sec:oportunidad}
Los sistemas de información hospitalarios actuales generan volúmenes de datos que hace apenas dos décadas habrían resultado inmanejables. Cada ingreso produce decenas o cientos de registros estructurados: analíticas seriadas, prescripciones farmacológicas, constantes vitales, diagnósticos codificados, procedimientos realizados. Sin embargo, la mayor parte de esta información se utiliza exclusivamente para el seguimiento individual del paciente y rara vez se explota de forma sistemática para la estratificación del riesgo o la detección temprana de complicaciones.

Las técnicas de aprendizaje automático ofrecen la posibilidad de extraer patrones predictivos de estos repositorios de datos clínicos. A diferencia de los modelos estadísticos tradicionales, los algoritmos de machine learning pueden manejar un número elevado de variables, capturar relaciones no lineales y adaptarse a la heterogeneidad propia de los datos hospitalarios reales. En el campo específico de la predicción de AKI, diversos estudios han demostrado que los modelos de aprendizaje automático pueden alcanzar un rendimiento discriminativo superior al de las escalas de riesgo convencionales \parencite{tomasev2019nature, churpek2020validation, rank2022mlaki}.

No obstante, la literatura también señala limitaciones importantes. Muchos de los modelos publicados presentan alto riesgo de sesgo metodológico, carecen de validación externa y se han desarrollado sobre poblaciones específicas (principalmente estadounidenses o del norte de Europa) que pueden no ser directamente extrapolables a otros contextos \parencite{shi2025aiaki}. Esta circunstancia subraya la necesidad de desarrollar y evaluar modelos en entornos clínicos diversos, incluyendo hospitales españoles con sus propias características poblacionales y organizativas.

\section{Contexto institucional}
\label{sec:institucional}
Este trabajo se desarrolla en el marco de HM Hospitales, con acceso a la base de datos clínica Doctoris. Esta circunstancia ofrece una oportunidad poco habitual en trabajos académicos de grado: la posibilidad de trabajar con datos hospitalarios reales, no simulados ni procedentes de repositorios públicos internacionales.

Trabajar con datos reales introduce complejidades adicionales frente a los datasets de laboratorio. La información clínica presenta valores ausentes, codificaciones heterogéneas, registros incompletos y particularidades derivadas de la práctica asistencial cotidiana. Sin embargo, esta misma complejidad confiere al trabajo una relevancia aplicada que no tendría si se limitara a reproducir experimentos sobre bases de datos previamente depuradas.

El acceso a Doctoris permite además construir una cohorte con características propias del sistema sanitario español, contribuyendo a llenar un vacío en la literatura, dominada por estudios realizados en otros países con sistemas de salud y perfiles poblacionales diferentes.

\section{Pertinencia del enfoque adoptado}
\label{sec:enfoque}
La elección de modelos de aprendizaje automático interpretables (regresión logística y Random Forest) responde a un criterio deliberado. Aunque existen arquitecturas más complejas, como las redes neuronales profundas utilizadas por Tomašev et al. \parencite{tomasev2019nature}, estas requieren volúmenes de datos considerablemente mayores y plantean retos de interpretabilidad que dificultan su aceptación en entornos clínicos. Para una prueba de concepto desarrollada en el contexto de un TFG, con las limitaciones de tamaño muestral inherentes a una cohorte monocéntrica, resulta más adecuado emplear modelos cuyo funcionamiento interno pueda analizarse y cuyas predicciones puedan explicarse en términos clínicamente comprensibles.

La inclusión de variables farmacológicas como predictores merece también justificación específica. Diversos fármacos de uso hospitalario habitual presentan potencial nefrotóxico bien documentado: los antiinflamatorios no esteroideos alteran la hemodinámica glomerular, los aminoglucósidos y la vancomicina ejercen toxicidad tubular directa, y los inhibidores del sistema renina-angiotensina pueden precipitar deterioro funcional en pacientes con perfusión renal comprometida \parencite{perazella2019drug, perazella2022drugaki}. La exposición a estos agentes constituye un factor de riesgo modificable, lo que confiere especial interés a su incorporación en modelos predictivos orientados a la prevención.

\section{Alcance y limitaciones asumidas}
\label{sec:alcance}
Resulta importante explicitar desde el inicio que este trabajo constituye una prueba de concepto académica, no una herramienta lista para despliegue clínico. El objetivo no es desarrollar un sistema de alertas que pueda implantarse mañana en un hospital, sino demostrar la viabilidad técnica del enfoque, evaluar el rendimiento de modelos baseline sobre datos reales e identificar qué variables parecen más relevantes para la predicción.

Las limitaciones propias de un TFG: tamaño muestral moderado, ausencia de validación externa, restricciones temporales para explorar múltiples arquitecturas. No invalidan el interés del trabajo, sino que definen su alcance. Si los resultados son prometedores, el siguiente paso natural sería ampliar la cohorte, incorporar nuevas variables y validar el modelo en otros centros o periodos temporales. Pero esas extensiones quedan, por definición, fuera del ámbito del presente trabajo de grado.

En síntesis, la justificación de este TFG descansa sobre tres pilares: la relevancia clínica de la AKI como problema prevenible con alto impacto pronóstico, la oportunidad tecnológica que ofrecen los datos hospitalarios estructurados y las técnicas de aprendizaje automático, y la disponibilidad de un entorno real (HM Hospitales, base de datos Doctoris) donde poner a prueba el enfoque propuesto. La confluencia de estos factores hace del problema abordado un campo de estudio con interés tanto académico como potencialmente aplicado.

%% --- CAPÍTULO 4: OBJETIVOS ---
\chapter{Objetivos}
\label{chap:objetivos}

\section{Objetivo general}
\label{sec:obj_general}

Desarrollar y evaluar un modelo de aprendizaje automático para predecir la aparición de lesión renal aguda en las primeras 48 horas de ingreso hospitalario, utilizando datos clínicos disponibles al momento del ingreso.

\section{Objetivos específicos}
\label{sec:obj_especificos}

\begin{enumerate}
    \item \textbf{OE1:} Identificar las variables clínicas, analíticas y farmacológicas disponibles en la base de datos hospitalaria que puedan estar asociadas al desarrollo de AKI.
    \item \textbf{OE2:} Construir un dataset estructurado mediante un proceso de extracción, limpieza y transformación de los datos clínicos.
    \item \textbf{OE3:} Realizar un análisis exploratorio para caracterizar la cohorte y detectar patrones relacionados con el desarrollo de AKI.
    \item \textbf{OE4:} Entrenar y comparar modelos de regresión logística y Random Forest para la predicción de AKI, evaluando su rendimiento mediante métricas apropiadas para problemas desbalanceados.
    \item \textbf{OE5:} Analizar la importancia de las variables predictoras e interpretar los resultados desde una perspectiva clínica.
\end{enumerate}

%% --- CAPÍTULO 5: METODOLOGÍA ---
\chapter{Metodología}
\label{chap:metodologia}

Este capítulo describe el enfoque metodológico seguido para el desarrollo del trabajo. Se presenta el entorno técnico utilizado, el proceso de definición de la cohorte de estudio, la construcción de la variable objetivo, las variables predictoras seleccionadas, el preprocesamiento aplicado y la estrategia de modelado adoptada.

\section{Entorno técnico}

El desarrollo del proyecto requirió la combinación de varias herramientas orientadas a la extracción, preparación, análisis y modelado de datos clínicos. La selección de estas herramientas respondió a criterios de disponibilidad institucional, robustez, adecuación metodológica y capacidad para construir un flujo de trabajo reproducible.

\subsection{Base de datos y extracción: SQL Server y Doctoris}

La fuente primaria de datos fue la base de datos clínica Doctoris, utilizada en los hospitales del grupo HM para el registro estructurado de información relacionada con pacientes, ingresos, resultados de laboratorio y tratamientos farmacológicos. A diferencia de otros enfoques basados en texto libre o historias clínicas no estructuradas, este trabajo utilizó exclusivamente datos tabulares y codificados, lo que permitió desarrollar un proceso de extracción más controlado y reproducible.

El acceso a Doctoris se realizó mediante Microsoft SQL Server. Esta plataforma permitió construir consultas complejas para definir la cohorte de estudio, integrar información procedente de múltiples tablas (pacientes, ingresos, laboratorio, medicación) y generar un dataset estructurado con una fila por ingreso hospitalario. La lógica de extracción se implementó mediante una query maestra que utilizó expresiones comunes de tabla (CTEs), funciones de ventana y agregaciones condicionales para transformar los datos clínicos brutos en una estructura analítica coherente.

Las principales tablas utilizadas fueron:

\begin{itemize}
    \item \textbf{Pacientes}: información demográfica básica, incluyendo identificador interno, identificador anonimizado, fecha de nacimiento y sexo.
    \item \textbf{Ingresos}: datos del episodio hospitalario, incluyendo identificador de ingreso, fechas de ingreso y alta, hospital y servicio.
    \item \textbf{Laboratorio\_Orden} y \textbf{LaboratorioResultados\_lab}: resultados analíticos con código de prueba, valor numérico, unidades y fecha de extracción.
    \item \textbf{TratamientoFarm} y tablas auxiliares de farmacología: prescripciones farmacológicas durante el ingreso, con información del principio activo.
\end{itemize}

El resultado de esta fase fue un fichero CSV con 1.175 filas (una por ingreso) y 24 columnas, que constituyó el punto de partida para el análisis posterior en Python.

\subsection{Entorno de análisis: Python y Jupyter Notebook}

El núcleo del trabajo analítico se desarrolló en Python, utilizando Jupyter Notebook como entorno de trabajo interactivo. Python fue seleccionado por su amplia adopción en ciencia de datos, su versatilidad para el tratamiento de datos estructurados y su ecosistema consolidado de librerías para aprendizaje automático.

La estructura del proyecto se organizó en varios notebooks especializados, cada uno centrado en una fase concreta del pipeline:

\begin{enumerate}
    \item Exploración inicial del dato
    \item Limpieza y preprocesamiento
    \item Ingeniería de variables
    \item Construcción del dataset final
    \item Entrenamiento de modelos
    \item Evaluación de modelos
    \item Interpretación de modelos
    \item Demostración sobre casos simulados
\end{enumerate}

Esta organización modular permitió separar claramente las diferentes etapas del trabajo, facilitando la trazabilidad de las decisiones y la reproducibilidad del análisis.

\subsection{Librerías utilizadas}

Para el desarrollo del pipeline analítico se utilizaron las siguientes librerías del ecosistema científico de Python:

\textbf{Pandas} fue la librería principal para la manipulación de datos tabulares. Se utilizó en tareas como la carga del fichero CSV, la asignación de tipos de datos, la detección y tratamiento de valores nulos, la transformación de variables y la construcción de los diferentes datasets intermedios.

\textbf{NumPy} se empleó como apoyo para operaciones numéricas y transformaciones sobre arrays, complementando las funcionalidades de Pandas en ciertas operaciones de cálculo.

\textbf{Matplotlib} y \textbf{Seaborn} se utilizaron para la representación gráfica de resultados. Estas librerías permitieron visualizar distribuciones de variables, matrices de correlación, curvas ROC y otros elementos necesarios para comprender el comportamiento del dataset y de los modelos entrenados.

\textbf{Scikit-learn} constituyó la librería central para el desarrollo de los modelos de aprendizaje automático. Se empleó para la separación entre conjuntos de entrenamiento y prueba, la imputación de valores ausentes, el entrenamiento de los modelos (regresión logística y Random Forest), el cálculo de métricas de rendimiento y la construcción de pipelines de preprocesamiento.

\textbf{Joblib} se utilizó para la serialización y almacenamiento de los modelos entrenados, permitiendo su reutilización posterior en la fase de demostración.

\subsection{Control de versiones}

Como práctica de buena gestión del proyecto, se utilizó Git como sistema de control de versiones. Esto permitió mantener un historial de cambios en el código, facilitar la recuperación de versiones anteriores y documentar la evolución del trabajo a lo largo del tiempo. Aunque el proyecto no involucró desarrollo colaborativo, el uso de control de versiones constituye una práctica recomendada en proyectos de ciencia de datos que favorece la reproducibilidad y la trazabilidad.

\section{Definición de la cohorte}

La cohorte del estudio se definió con el objetivo de obtener un conjunto de ingresos hospitalarios que cumplieran criterios específicos de inclusión y que dispusieran de la información necesaria para construir tanto la variable objetivo como las variables predictoras.

\subsection{Unidad de análisis}

La unidad de análisis seleccionada fue el ingreso hospitalario, no el paciente. Esta decisión tiene sentido clínico: un mismo paciente puede tener riesgo muy diferente de desarrollar AKI en distintas hospitalizaciones, dependiendo de su situación basal, los tratamientos recibidos y el contexto clínico de cada episodio. Tratar cada ingreso como una observación independiente permite capturar esta variabilidad.

\subsection{Criterios de inclusión}

Para formar parte de la cohorte, cada ingreso debía cumplir los siguientes criterios:

\begin{enumerate}
    \item \textbf{Fecha de ingreso válida}: se excluyeron registros con fechas ausentes o inconsistentes.
    \item \textbf{Periodo de estudio}: se restringió el análisis a ingresos a partir de enero de 2024, periodo con mayor densidad de datos clínicos enlazables.
    \item \textbf{Identificador anonimizado disponible}: necesario para mantener la trazabilidad del paciente dentro del marco de anonimización.
    \item \textbf{Analítica disponible durante el ingreso}: imprescindible para poder definir tanto las variables predictoras analíticas como la variable objetivo basada en creatinina.
\end{enumerate}

\subsection{Proceso de enlace entre ingresos y laboratorio}

Un aspecto técnicamente relevante fue que la relación entre ingresos y resultados de laboratorio no venía dada de forma directa en la base de datos. Fue necesario reconstruir el vínculo mediante la coincidencia del identificador de paciente y la verificación de que la fecha de la muestra analítica estuviera comprendida entre la fecha de ingreso y la fecha de alta del episodio.

Esta lógica de enlace temporal permitió asociar los resultados de laboratorio al contexto clínico del ingreso correspondiente, aunque también introdujo una restricción: solo pudieron incluirse en la cohorte aquellos episodios con laboratorio efectivamente enlazable dentro de dicha ventana temporal.

\subsection{Tamaño final de la cohorte}

Tras aplicar todos los criterios de inclusión, la cohorte final quedó compuesta por \textbf{1.175 ingresos hospitalarios}. Esta cifra representa una fracción reducida del volumen bruto de ingresos históricos disponibles en la base de datos, lo que refleja una realidad común en proyectos de inteligencia artificial aplicada a datos clínicos: la disponibilidad efectiva de información utilizable para modelado suele ser muy inferior al tamaño total de las bases de datos.

\section{Construcción de la variable objetivo}

La variable objetivo del estudio fue \textbf{AKI\_48h}, definida como la aparición de lesión renal aguda en las 48 horas posteriores al valor basal de creatinina del ingreso.

\subsection{Definición operativa}

El criterio utilizado para definir AKI fue un incremento de la creatinina sérica de al menos 0,3 mg/dL respecto al valor basal, siguiendo una adaptación simplificada del criterio KDIGO de incremento absoluto en 48 horas \parencite{kdigo2012aki}.

Formalmente, la variable se definió como:

\begin{equation}
\text{AKI\_48h} = 
\begin{cases}
1 & \text{si } \text{MaxCreatinina48h} \geq \text{BaselineCreatinina} + 0{,}3 \\
0 & \text{en caso contrario}
\end{cases}
\end{equation}

Donde:
\begin{itemize}
    \item \textbf{BaselineCreatinina}: primera creatinina disponible del ingreso, utilizada como valor de referencia.
    \item \textbf{MaxCreatinina48h}: valor máximo de creatinina observado en las 48 horas posteriores al baseline.
\end{itemize}

\subsection{Proceso de construcción}

La construcción de la variable objetivo siguió varios pasos:

\begin{enumerate}
    \item Identificación de todas las determinaciones de creatinina (código de prueba CREA) realizadas durante cada ingreso.
    \item Ordenación cronológica de las determinaciones dentro de cada episodio.
    \item Selección de la primera creatinina como valor baseline.
    \item Cálculo del valor máximo de creatinina en las 48 horas posteriores al baseline.
    \item Aplicación del criterio de incremento $\geq$ 0,3 mg/dL para etiquetar cada ingreso.
\end{enumerate}

\subsection{Distribución de la variable objetivo}

La distribución resultante mostró un marcado desbalanceo de clases:

\begin{itemize}
    \item \textbf{Casos positivos (AKI\_48h = 1)}: 29 ingresos (2,5\%)
    \item \textbf{Casos negativos (AKI\_48h = 0)}: 1.146 ingresos (97,5\%)
\end{itemize}

Esta prevalencia tan baja del evento positivo condicionó de forma directa tanto la estrategia de modelado como la selección de métricas de evaluación, aspectos que se desarrollan más adelante en este capítulo.

\section{Variables predictoras}

El conjunto de variables predictoras se seleccionó buscando un equilibrio entre relevancia clínica, disponibilidad en la base de datos y viabilidad metodológica. Las variables se organizaron en cuatro categorías: demográficas, analíticas, farmacológicas y derivadas.

\subsection{Variables demográficas}

Se incluyeron dos variables demográficas básicas:

\begin{itemize}
    \item \textbf{Sexo}: variable binaria (0 = mujer, 1 = hombre), recodificada a partir de la codificación original de la base de datos.
    \item \textbf{EdadAproxIngreso}: edad del paciente en años al momento del ingreso, calculada como la diferencia entre la fecha de nacimiento y la fecha de ingreso.
\end{itemize}

La edad es un factor de riesgo bien establecido para el desarrollo de AKI, dado que la función renal tiende a deteriorarse con el envejecimiento y los pacientes mayores suelen presentar más comorbilidades y mayor exposición a fármacos potencialmente nefrotóxicos.

\subsection{Variables analíticas}

Se seleccionaron tres variables analíticas correspondientes a la primera determinación disponible durante el ingreso:

\begin{itemize}
    \item \textbf{BaselineCreatinina}: creatinina sérica basal en mg/dL. Esta variable refleja el estado de la función renal al inicio del episodio y constituye el punto de referencia para la definición de AKI.
    \item \textbf{Potasio\_mmol\_L}: concentración de potasio sérico en mmol/L. Las alteraciones del potasio pueden asociarse a disfunción renal y a determinadas situaciones clínicas de riesgo.
    \item \textbf{Sodio\_mmol\_L}: concentración de sodio sérico en mmol/L. Las alteraciones del sodio, especialmente la hiponatremia, pueden indicar situaciones de sobrecarga de volumen o depleción que se relacionan con el riesgo renal.
\end{itemize}

La variable \textbf{Urea\_mg\_dL} fue inicialmente extraída pero posteriormente excluida debido a su muy baja disponibilidad (menos del 4\% de los ingresos tenían este dato), lo que la hacía inviable para su uso como predictor.

\subsection{Variables farmacológicas}

Se construyeron seis variables binarias indicando la exposición a grupos farmacológicos con potencial nefrotóxico durante el ingreso:

\begin{itemize}
    \item \textbf{AINE}: exposición a antiinflamatorios no esteroideos (ibuprofeno, naproxeno, diclofenaco, etc.).
    \item \textbf{IECA}: exposición a inhibidores de la enzima convertidora de angiotensina (enalapril, ramipril, etc.).
    \item \textbf{ARAII}: exposición a antagonistas de los receptores de angiotensina II (losartán, valsartán, etc.).
    \item \textbf{DIURETICO}: exposición a diuréticos (furosemida, torasemida, hidroclorotiazida, etc.).
    \item \textbf{VANCOMICINA}: exposición a vancomicina, antibiótico con nefrotoxicidad bien documentada.
    \item \textbf{AMINOGLUCOSIDO}: exposición a aminoglucósidos (gentamicina, amikacina, etc.), antibióticos con conocido potencial nefrotóxico.
\end{itemize}

Estas variables se construyeron a partir de la tabla de tratamiento farmacológico, enlazando con las tablas de catálogo de fármacos para identificar el principio activo y agregando a nivel de ingreso mediante búsquedas textuales sobre nombres de medicamentos.

La justificación clínica para incluir estos grupos farmacológicos se basa en la literatura sobre nefrotoxicidad inducida por fármacos \parencite{perazella2019drug, perazella2022drugaki}. Los AINEs alteran la hemodinámica glomerular al inhibir las prostaglandinas vasodilatadoras. Los IECAs y ARA-II pueden precipitar deterioro funcional en pacientes con perfusión renal comprometida. Los aminoglucósidos y la vancomicina ejercen toxicidad tubular directa. Los diuréticos, aunque no son directamente nefrotóxicos, pueden contribuir a situaciones de depleción de volumen que favorecen el daño renal.

\subsection{Variables derivadas}

Se crearon dos variables derivadas con el objetivo de enriquecer la representación del problema:

\begin{itemize}
    \item \textbf{Edad\_ge\_65}: variable binaria que indica si el paciente tiene 65 años o más. El umbral de 65 años es un punto de corte habitual en la literatura para identificar pacientes de edad avanzada con mayor vulnerabilidad.
    \item \textbf{N\_Farmacos\_Riesgo}: variable numérica que cuenta el número de grupos farmacológicos de riesgo a los que está expuesto el paciente (suma de AINE, IECA, ARAII, DIURETICO, VANCOMICINA y AMINOGLUCOSIDO). Esta variable sintetiza la carga farmacológica potencialmente nefrotóxica en una única medida.
\end{itemize}

\subsection{Variables excluidas}

Varias variables disponibles en la extracción inicial fueron excluidas del modelado por diferentes motivos:

\begin{itemize}
    \item \textbf{Identificadores} (IdIngreso, UUIdPaciente, IdHospital, IdDoctor): no constituyen predictores clínicos y su inclusión introduciría señales espurias.
    \item \textbf{Variables temporales} (FechaIngreso, FechaAlta, fechas de laboratorio): podrían utilizarse para construir nuevas características, pero se optó por un enfoque inicial más simple.
    \item \textbf{MaxCreatinina48h}: esta variable forma parte de la definición del desenlace, por lo que su inclusión como predictor constituiría una fuga de información.
    \item \textbf{Creatinina\_mg\_dL}: presentaba una correlación de 1,00 con BaselineCreatinina (coincidencia exacta en el 99,57\% de los casos), por lo que se eliminó por redundancia.
    \item \textbf{Urea\_mg\_dL}: disponibilidad inferior al 4\%, inviable para modelado.
\end{itemize}

\subsection{Resumen del conjunto de predictores}

El dataset final de modelado quedó compuesto por \textbf{13 variables predictoras}:

\begin{table}[h]
\centering
\begin{tabular}{lll}
\hline
\textbf{Variable} & \textbf{Tipo} & \textbf{Descripción} \\
\hline
Sexo & Binaria & 0 = mujer, 1 = hombre \\
EdadAproxIngreso & Continua & Edad en años \\
Edad\_ge\_65 & Binaria & 1 si edad $\geq$ 65 \\
BaselineCreatinina & Continua & Creatinina basal (mg/dL) \\
Potasio\_mmol\_L & Continua & Potasio sérico (mmol/L) \\
Sodio\_mmol\_L & Continua & Sodio sérico (mmol/L) \\
AINE & Binaria & Exposición a AINEs \\
IECA & Binaria & Exposición a IECAs \\
ARAII & Binaria & Exposición a ARA-II \\
DIURETICO & Binaria & Exposición a diuréticos \\
VANCOMICINA & Binaria & Exposición a vancomicina \\
AMINOGLUCOSIDO & Binaria & Exposición a aminoglucósidos \\
N\_Farmacos\_Riesgo & Discreta & Número de fármacos de riesgo (0-6) \\
\hline
\end{tabular}
\caption{Variables predictoras incluidas en el modelo}
\label{tab:variables}
\end{table}

\section{Preprocesamiento de datos}

Antes del modelado, los datos requirieron un proceso de limpieza y transformación para garantizar su calidad y adecuación a los algoritmos de aprendizaje automático.

\subsection{Validación estructural}

Se verificó que el dataset cumplía las características esperadas:

\begin{itemize}
    \item 1.175 observaciones correspondientes a ingresos únicos (sin duplicados).
    \item Granularidad correcta: una fila por ingreso hospitalario.
    \item Alineación correcta de columnas tras la importación del CSV.
\end{itemize}

\subsection{Conversión de tipos de datos}

Durante la importación se detectó que los valores decimales habían sido exportados con coma decimal en lugar de punto (formato europeo). Esto impedía su interpretación automática como valores numéricos. Se aplicó una transformación para reemplazar las comas por puntos y convertir las columnas afectadas a tipo numérico.

Las variables se tiparon de la siguiente manera:
\begin{itemize}
    \item Variables continuas (creatinina, potasio, sodio): tipo float.
    \item Variables binarias y discretas: tipo int.
    \item Variable objetivo (AKI\_48h): tipo int.
\end{itemize}

\subsection{Tratamiento de valores ausentes}

El análisis de completitud mostró que la mayoría de las variables tenían cobertura completa, con excepción de las tres variables analíticas continuas:

\begin{itemize}
    \item Potasio\_mmol\_L: 4,51\% de valores ausentes
    \item Sodio\_mmol\_L: 4,17\% de valores ausentes
    \item BaselineCreatinina: 0,17\% de valores ausentes
\end{itemize}

Dado que estos porcentajes eran relativamente bajos y que la eliminación de casos incompletos habría reducido aún más el ya limitado número de eventos positivos, se optó por una estrategia de \textbf{imputación por mediana}. Esta técnica es robusta frente a valores extremos y resulta apropiada para biomarcadores clínicos cuya distribución puede no ser simétrica.

La imputación se integró dentro del pipeline de modelado para garantizar que los valores de imputación se calcularan exclusivamente sobre el conjunto de entrenamiento y se aplicaran después al conjunto de prueba, evitando así fuga de información.

\subsection{Recodificación de variables}

La variable Sexo, originalmente codificada como 1/2 en la base de datos, fue transformada a codificación 0/1 para facilitar su interpretación en los modelos (0 = mujer, 1 = hombre).

\section{Estrategia de modelado}

La estrategia de modelado se diseñó teniendo en cuenta las características específicas del problema: tamaño muestral moderado, fuerte desbalanceo de clases y objetivo de construir modelos interpretables.

\subsection{Partición de datos}

El dataset se dividió en un conjunto de entrenamiento (80\%) y un conjunto de prueba (20\%), utilizando una semilla aleatoria fija para garantizar la reproducibilidad. La partición se realizó de forma \textbf{estratificada}, manteniendo la proporción de clases aproximadamente constante en ambos subconjuntos.

Tras la partición:
\begin{itemize}
    \item \textbf{Conjunto de entrenamiento}: 940 observaciones (917 negativos, 23 positivos)
    \item \textbf{Conjunto de prueba}: 235 observaciones (229 negativos, 6 positivos)
\end{itemize}

La estratificación resultó imprescindible dado el bajo número de eventos positivos. Una partición aleatoria simple podría haber generado subconjuntos con distribuciones de clase poco representativas.

\subsection{Modelos seleccionados}

Se entrenaron dos modelos de clasificación:

\textbf{Regresión logística}: modelo lineal que estima la probabilidad de pertenencia a la clase positiva mediante una función logística. Sus principales ventajas son la interpretabilidad (los coeficientes pueden expresarse como odds ratios) y su amplio uso en investigación biomédica. Se configuró con un máximo de 1.000 iteraciones para asegurar la convergencia.

\textbf{Random Forest}: modelo de ensamble basado en múltiples árboles de decisión entrenados sobre submuestras aleatorias de los datos. Cada árbol vota por una clase y la predicción final se obtiene por agregación. Sus ventajas incluyen la capacidad de capturar relaciones no lineales, la robustez frente al sobreajuste y la posibilidad de estimar la importancia relativa de las variables. Se configuró con 200 árboles (n\_estimators = 200).

Ambos modelos se configuraron con \textbf{class\_weight = 'balanced'}, opción que pondera automáticamente las clases de forma inversamente proporcional a su frecuencia. Esto ayuda a contrarrestar la tendencia de los algoritmos a favorecer la clase mayoritaria en problemas desbalanceados.

\subsection{Pipeline de modelado}

El preprocesamiento y el modelado se integraron en un pipeline único utilizando las herramientas de Scikit-learn. Esto garantizó que:

\begin{enumerate}
    \item Los parámetros de imputación se ajustaran solo sobre el conjunto de entrenamiento.
    \item El mismo preprocesamiento se aplicara de forma consistente al conjunto de prueba.
    \item El flujo fuera reproducible y menos propenso a errores.
\end{enumerate}

\subsection{Métricas de evaluación}

Dado el fuerte desbalanceo de clases, la \textit{accuracy} (exactitud global) resultaría engañosa como métrica principal: un modelo que predijera siempre la clase negativa alcanzaría un 97,5\% de accuracy sin aportar ninguna utilidad práctica.

Por este motivo, la evaluación se basó en métricas más apropiadas para problemas con evento infrecuente:

\begin{itemize}
    \item \textbf{ROC-AUC}: área bajo la curva ROC, que mide la capacidad global del modelo para discriminar entre clases independientemente del umbral de decisión.
    \item \textbf{Recall} (sensibilidad): proporción de casos positivos correctamente identificados. Especialmente relevante en contextos clínicos donde interesa detectar a los pacientes en riesgo.
    \item \textbf{Precision}: proporción de predicciones positivas que son correctas. Indica la fiabilidad de las alertas generadas.
    \item \textbf{F1-score}: media armónica de precision y recall, útil como medida de equilibrio entre ambas.
    \item \textbf{Matriz de confusión}: permite visualizar los verdaderos positivos, verdaderos negativos, falsos positivos y falsos negativos.
\end{itemize}

Adicionalmente, se realizó un análisis del umbral de decisión para evaluar cómo variaba el rendimiento del modelo al modificar el punto de corte para la clasificación binaria, dado que el umbral estándar de 0,5 puede no ser óptimo en problemas muy desbalanceados.

%% --- CAPÍTULO 6: DESARROLLO ---
\chapter{Desarrollo del Trabajo}
\label{chap:desarrollo}

Este capítulo describe el desarrollo técnico del proyecto, desde la exploración inicial de los datos hasta el entrenamiento de los modelos predictivos. Se detallan las decisiones tomadas en cada fase, los problemas encontrados y las soluciones adoptadas.

\section{Organización del pipeline de análisis}

El trabajo técnico se estructuró en una secuencia de notebooks de Jupyter, cada uno dedicado a una fase específica del proceso analítico. Esta organización modular permitió mantener la trazabilidad de las decisiones, facilitar la depuración de errores y garantizar la reproducibilidad del análisis.

Los notebooks desarrollados fueron los siguientes:

\begin{enumerate}
    \item \textbf{01\_data\_exploration}: exploración inicial del dataset importado, validación estructural y análisis descriptivo preliminar.
    \item \textbf{02\_data\_cleaning}: limpieza de datos, corrección de tipos, tratamiento de incidencias de importación.
    \item \textbf{03\_feature\_engineering}: análisis de correlaciones, detección de redundancias, creación de variables derivadas.
    \item \textbf{04\_dataset\_creation}: construcción del dataset final de modelado, selección definitiva de predictores.
    \item \textbf{05\_model\_training}: partición de datos, configuración de pipelines, entrenamiento de modelos.
    \item \textbf{06\_model\_evaluation}: cálculo de métricas, matrices de confusión, curvas ROC, comparación de modelos.
    \item \textbf{07\_model\_interpretation}: análisis de coeficientes e importancias, interpretación clínica de resultados.
    \item \textbf{08\_demo\_prediction}: demostración del sistema sobre casos clínicos simulados.
\end{enumerate}

Esta estructura refleja un flujo de trabajo típico en proyectos de ciencia de datos aplicada, donde cada fase construye sobre los resultados de la anterior y produce artefactos (datasets intermedios, modelos serializados, figuras) que quedan documentados y pueden reutilizarse.

\section{Análisis exploratorio inicial}

La primera fase del trabajo analítico consistió en examinar el dataset exportado desde SQL Server para validar su estructura, caracterizar las variables disponibles e identificar posibles problemas que pudieran comprometer las fases posteriores.

\subsection{Validación de la estructura del dataset}

Tras la importación del fichero CSV, se realizaron las siguientes comprobaciones:

\begin{itemize}
    \item \textbf{Dimensiones}: el dataset contenía 1.175 filas y 24 columnas, coincidiendo con lo esperado según la query de extracción.
    \item \textbf{Unicidad}: se verificó que no existían filas duplicadas y que cada valor de IdIngreso era único, confirmando la granularidad de una observación por ingreso hospitalario.
    \item \textbf{Alineación de columnas}: se comprobó que los nombres de las columnas correspondían correctamente con su contenido, descartando desplazamientos producidos durante la exportación.
\end{itemize}

Durante esta validación se detectó una incidencia relevante: el fichero CSV utilizaba punto y coma como separador de campos (formato europeo) y no incluía fila de cabecera. Fue necesario especificar explícitamente el separador en la importación y asignar manualmente los nombres de las columnas según el orden definido en la query SQL.

\subsection{Análisis de completitud}

El estudio de valores ausentes reveló un patrón heterogéneo de disponibilidad entre variables:

\begin{table}[h]
\centering
\begin{tabular}{lrr}
\hline
\textbf{Variable} & \textbf{Valores ausentes} & \textbf{Porcentaje} \\
\hline
Urea\_mg\_dL & 1.134 & 96,5\% \\
Potasio\_mmol\_L & 53 & 4,5\% \\
Sodio\_mmol\_L & 49 & 4,2\% \\
BaselineCreatinina & 2 & 0,2\% \\
Resto de variables & 0 & 0,0\% \\
\hline
\end{tabular}
\caption{Análisis de valores ausentes en el dataset}
\label{tab:missing}
\end{table}

La variable Urea\_mg\_dL presentaba una disponibilidad extremadamente baja, con menos del 4\% de los ingresos con valor registrado. Esta situación la hacía inviable como predictor, por lo que se decidió su exclusión del análisis. Las variables de potasio y sodio mostraban tasas de ausencia moderadas pero manejables mediante imputación. La creatinina basal tenía una completitud prácticamente total.

\subsection{Distribución de la variable objetivo}

El análisis de la variable AKI\_48h confirmó el marcado desbalanceo de clases anticipado:

\begin{itemize}
    \item Casos negativos (AKI\_48h = 0): 1.146 ingresos (97,5\%)
    \item Casos positivos (AKI\_48h = 1): 29 ingresos (2,5\%)
\end{itemize}

Esta distribución tiene implicaciones metodológicas importantes. Con solo 29 eventos positivos en toda la cohorte, y aproximadamente 6 en el conjunto de prueba tras la partición, las estimaciones de rendimiento están sujetas a considerable variabilidad. Un solo caso positivo correctamente clasificado o no puede hacer variar el recall en más de 15 puntos porcentuales.

El desbalanceo también condiciona la interpretación de las métricas. Un modelo que predijera siempre la clase negativa alcanzaría una exactitud del 97,5\%, pero sería completamente inútil para el propósito clínico de identificar pacientes en riesgo. Por este motivo, desde el inicio se planificó basar la evaluación en métricas como ROC-AUC, recall y F1-score, más informativas en este contexto.

\subsection{Incidencia del formato decimal}

Durante la exploración de las variables numéricas se detectó un problema técnico que habría comprometido el análisis de no haberse identificado a tiempo. Al intentar convertir las columnas analíticas (creatinina, potasio, sodio) a tipo numérico, se generaba un número anormalmente alto de valores nulos.

La inspección manual de los valores reveló la causa: los decimales habían sido exportados con coma en lugar de punto (por ejemplo, ``0,96'' en lugar de ``0.96''). Este formato, habitual en configuraciones regionales europeas, no es reconocido automáticamente por las funciones de conversión numérica de Python.

La solución consistió en aplicar una transformación previa que reemplazara las comas por puntos antes de la conversión a tipo float. Este paso, aparentemente menor, fue crítico para preservar la integridad de los biomarcadores analíticos, que constituyen predictores centrales del modelo.

\subsection{Estadísticos descriptivos de las variables continuas}

Una vez resueltos los problemas de formato, se calcularon los estadísticos descriptivos de las principales variables numéricas:

\begin{table}[h]
\centering
\begin{tabular}{lrrrrr}
\hline
\textbf{Variable} & \textbf{Media} & \textbf{DE} & \textbf{Mín} & \textbf{Mediana} & \textbf{Máx} \\
\hline
EdadAproxIngreso & 58,3 & 19,2 & 0 & 61 & 99 \\
BaselineCreatinina & 0,91 & 0,48 & 0,20 & 0,82 & 7,85 \\
Potasio\_mmol\_L & 4,21 & 0,52 & 2,60 & 4,20 & 7,10 \\
Sodio\_mmol\_L & 139,8 & 3,8 & 121 & 140 & 152 \\
\hline
\end{tabular}
\caption{Estadísticos descriptivos de las variables continuas}
\label{tab:descriptivos}
\end{table}

Los valores observados resultaron coherentes con lo esperado en una población hospitalaria general. La creatinina basal mostró una mediana de 0,82 mg/dL, con valores que oscilaban entre 0,20 y 7,85 mg/dL. Los valores extremos superiores corresponden probablemente a pacientes con enfermedad renal crónica previa, y no fueron eliminados al no constituir errores de medición sino situaciones clínicas reales.

\subsection{Distribución de las variables farmacológicas}

Las variables binarias de exposición farmacológica mostraron frecuencias variables:

\begin{table}[h]
\centering
\begin{tabular}{lr}
\hline
\textbf{Variable} & \textbf{Frecuencia (\%)} \\
\hline
DIURETICO & 18,7\% \\
AINE & 15,2\% \\
IECA & 8,3\% \\
ARAII & 6,1\% \\
VANCOMICINA & 2,8\% \\
AMINOGLUCOSIDO & 0,9\% \\
\hline
\end{tabular}
\caption{Frecuencia de exposición a fármacos de riesgo}
\label{tab:farmacos}
\end{table}

Los diuréticos y los AINEs fueron los grupos más frecuentes, lo cual tiene sentido dado su amplio uso hospitalario. La vancomicina y especialmente los aminoglucósidos mostraron frecuencias muy bajas, lo que limita la capacidad de los modelos para estimar de forma estable su contribución al riesgo.

\section{Ingeniería de variables}

La fase de ingeniería de variables se planteó desde una perspectiva deliberadamente conservadora. En lugar de generar un gran número de características derivadas, se priorizó la parsimonia y la interpretabilidad, teniendo en cuenta las limitaciones impuestas por el tamaño muestral y el número reducido de eventos positivos.

\subsection{Análisis de correlaciones}

Como paso previo a cualquier transformación, se analizó la matriz de correlaciones entre las variables numéricas del dataset. El hallazgo más relevante fue la correlación prácticamente perfecta (r = 1,00) entre las variables Creatinina\_mg\_dL y BaselineCreatinina.

Una inspección más detallada reveló que ambas variables coincidían exactamente en 1.170 de las 1.175 observaciones (99,57\%). Las cinco discrepancias correspondían a casos donde existían múltiples determinaciones de creatinina en fechas muy próximas, generando pequeñas diferencias entre el valor del primer laboratorio y el valor asignado como baseline.

Esta redundancia tenía implicaciones claras para el modelado:

\begin{itemize}
    \item Mantener ambas variables no aportaría información adicional al modelo.
    \item En regresión logística, la colinealidad perfecta generaría inestabilidad en la estimación de coeficientes.
    \item La interpretación de los resultados se complicaría innecesariamente.
\end{itemize}

Por estos motivos, se decidió conservar únicamente BaselineCreatinina, por ser la representación más coherente con la lógica temporal del estudio (valor de referencia para la definición del desenlace).

\subsection{Creación de variables derivadas}

Se generaron dos variables derivadas con justificación clínica clara:

\textbf{Edad\_ge\_65}: variable binaria que toma valor 1 si el paciente tiene 65 años o más, y 0 en caso contrario. El umbral de 65 años es un punto de corte ampliamente utilizado en la literatura médica para identificar pacientes de edad avanzada con mayor vulnerabilidad. Esta variable complementa a la edad continua, permitiendo capturar un posible efecto umbral que un modelo lineal no detectaría automáticamente.

\textbf{N\_Farmacos\_Riesgo}: variable discreta que cuenta el número de grupos farmacológicos de riesgo a los que está expuesto el paciente durante el ingreso. Se calcula como la suma de las seis variables farmacológicas binarias (AINE, IECA, ARAII, DIURETICO, VANCOMICINA, AMINOGLUCOSIDO), pudiendo tomar valores de 0 a 6.

La inclusión de N\_Farmacos\_Riesgo responde a la hipótesis de que la acumulación de exposiciones nefrotóxicas podría tener un efecto aditivo sobre el riesgo. En lugar de depender exclusivamente de la interpretación individual de cada fármaco, esta variable sintetiza la carga farmacológica global en una única medida.

\subsection{Relación preliminar con el desenlace}

Se exploró de forma descriptiva la relación entre las variables derivadas y la variable objetivo:

\begin{table}[h]
\centering
\begin{tabular}{lrr}
\hline
\textbf{Edad\_ge\_65} & \textbf{N} & \textbf{\% AKI} \\
\hline
0 (< 65 años) & 512 & 1,8\% \\
1 ($\geq$ 65 años) & 663 & 3,0\% \\
\hline
\end{tabular}
\caption{Distribución de AKI según grupo de edad}
\label{tab:edad_aki}
\end{table}

\begin{table}[h]
\centering
\begin{tabular}{lrr}
\hline
\textbf{N\_Farmacos\_Riesgo} & \textbf{N} & \textbf{\% AKI} \\
\hline
0 & 687 & 1,9\% \\
1 & 358 & 2,5\% \\
2 & 108 & 4,6\% \\
3+ & 22 & 9,1\% \\
\hline
\end{tabular}
\caption{Distribución de AKI según número de fármacos de riesgo}
\label{tab:nfarmacos_aki}
\end{table}

Los datos sugieren una tendencia hacia mayor frecuencia de AKI en pacientes de edad avanzada y en aquellos expuestos a múltiples fármacos de riesgo. Sin embargo, dado el bajo número de eventos positivos, estas diferencias deben interpretarse con cautela y no permiten establecer conclusiones inferenciales sólidas. Su utilidad principal fue confirmar que las variables derivadas mostraban un comportamiento coherente con las hipótesis clínicas subyacentes.

\subsection{Variables excluidas del modelado}

Durante esta fase se formalizó la exclusión de varias variables que no debían formar parte del conjunto de predictores:

\begin{itemize}
    \item \textbf{Identificadores} (IdIngreso, UUIdPaciente, IdHospital, IdSociedad, IdDoctor): campos administrativos sin valor predictivo clínico.
    \item \textbf{Fechas} (FechaIngreso, FechaAlta, PrimeraFechaLab, BaselineFecha): podrían usarse para derivar características temporales, pero se optó por un enfoque inicial más simple.
    \item \textbf{MaxCreatinina48h}: forma parte de la definición del desenlace; su inclusión constituiría fuga de información.
    \item \textbf{Creatinina\_mg\_dL}: redundante con BaselineCreatinina (correlación = 1,00).
    \item \textbf{Urea\_mg\_dL}: disponibilidad inferior al 4\%, inviable para modelado.
    \item \textbf{Diff\_Creatininas}: variable auxiliar creada para análisis, sin sentido como predictor clínico.
\end{itemize}

\section{Construcción del dataset final}

Tras el análisis exploratorio y la ingeniería de variables, se procedió a consolidar el dataset definitivo que serviría de entrada para los modelos predictivos.

\subsection{Selección final de variables}

El conjunto de predictores quedó definido por las siguientes 13 variables:

\begin{itemize}
    \item Variables demográficas: Sexo, EdadAproxIngreso, Edad\_ge\_65
    \item Variables analíticas: BaselineCreatinina, Potasio\_mmol\_L, Sodio\_mmol\_L
    \item Variables farmacológicas: AINE, IECA, ARAII, DIURETICO, VANCOMICINA, AMINOGLUCOSIDO
    \item Variables derivadas: N\_Farmacos\_Riesgo
\end{itemize}

La variable objetivo fue AKI\_48h.

\subsection{Estructura del dataset final}

El dataset final presentaba las siguientes características:

\begin{itemize}
    \item \textbf{Observaciones}: 1.175 ingresos hospitalarios
    \item \textbf{Predictores}: 13 variables
    \item \textbf{Variable objetivo}: 1 variable binaria
    \item \textbf{Valores ausentes}: concentrados en Potasio (4,5\%), Sodio (4,2\%) y BaselineCreatinina (0,2\%)
\end{itemize}

\subsection{Separación de predictores y objetivo}

Se realizó la separación formal entre la matriz de predictores X (1.175 × 13) y el vector objetivo y (1.175 × 1), siguiendo el formato estándar requerido por las librerías de aprendizaje automático de Python.

El dataset final se almacenó como fichero CSV independiente en la carpeta de datos procesados, permitiendo su reutilización directa en los notebooks posteriores sin necesidad de repetir todo el preprocesamiento.

\section{Entrenamiento de modelos}

La fase de entrenamiento se diseñó para producir modelos baseline sólidos y comparables, sin pretender una optimización exhaustiva de hiperparámetros.

\subsection{Partición de datos}

El dataset se dividió en conjunto de entrenamiento (80\%) y conjunto de prueba (20\%) mediante la función train\_test\_split de Scikit-learn, con los siguientes parámetros:

\begin{itemize}
    \item \texttt{test\_size = 0.2}
    \item \texttt{random\_state = 42} (semilla fija para reproducibilidad)
    \item \texttt{stratify = y} (partición estratificada por la variable objetivo)
\end{itemize}

La estratificación garantizó que la proporción de casos positivos fuera aproximadamente igual en ambos subconjuntos:

\begin{itemize}
    \item Entrenamiento: 940 observaciones (917 negativos, 23 positivos)
    \item Prueba: 235 observaciones (229 negativos, 6 positivos)
\end{itemize}

\subsection{Construcción del pipeline de preprocesamiento}

Para garantizar que el preprocesamiento se aplicara de forma consistente y sin fuga de información, se construyó un pipeline integrado que incluía:

\begin{enumerate}
    \item \textbf{Imputación}: se utilizó SimpleImputer con estrategia de mediana para las variables continuas con valores ausentes. La mediana se calculó exclusivamente sobre el conjunto de entrenamiento.
    \item \textbf{Modelado}: el modelo de clasificación (regresión logística o Random Forest) se añadió como paso final del pipeline.
\end{enumerate}

Este enfoque aseguró que:
\begin{itemize}
    \item Los valores de imputación no incorporaran información del conjunto de prueba.
    \item El mismo preprocesamiento se aplicara automáticamente a nuevos datos.
    \item El flujo fuera reproducible y menos propenso a errores manuales.
\end{itemize}

\subsection{Configuración de la regresión logística}

El modelo de regresión logística se configuró con los siguientes parámetros:

\begin{itemize}
    \item \texttt{max\_iter = 1000}: número máximo de iteraciones para garantizar la convergencia del algoritmo de optimización.
    \item \texttt{class\_weight = 'balanced'}: ponderación automática de las clases inversamente proporcional a su frecuencia, para contrarrestar el desbalanceo.
    \item \texttt{random\_state = 42}: semilla fija para reproducibilidad.
    \item \texttt{solver = 'lbfgs'}: algoritmo de optimización por defecto, adecuado para datasets de tamaño moderado.
\end{itemize}

La opción de ponderación balanceada hace que el algoritmo penalice más los errores sobre la clase minoritaria, incentivando al modelo a no ignorarla completamente en favor de maximizar la exactitud global.

\subsection{Configuración del Random Forest}

El modelo Random Forest se configuró con los siguientes parámetros:

\begin{itemize}
    \item \texttt{n\_estimators = 200}: número de árboles en el bosque.
    \item \texttt{class\_weight = 'balanced'}: misma estrategia de ponderación que en la regresión logística.
    \item \texttt{random\_state = 42}: semilla fija para reproducibilidad.
    \item Resto de parámetros: valores por defecto de Scikit-learn.
\end{itemize}

No se realizó una búsqueda exhaustiva de hiperparámetros (grid search o random search) porque el objetivo era establecer una línea base razonable, no optimizar el rendimiento al máximo. Con solo 23 eventos positivos en el conjunto de entrenamiento, una optimización agresiva habría incrementado el riesgo de sobreajuste sin garantías de mejora real en datos nuevos.

\subsection{Proceso de entrenamiento}

El entrenamiento de ambos modelos se realizó mediante el método \texttt{fit()} sobre el conjunto de entrenamiento. El proceso fue rápido (segundos) dado el tamaño moderado del dataset.

Tras el entrenamiento, los pipelines completos (preprocesamiento + modelo) se serializaron mediante Joblib para su reutilización posterior en la fase de demostración.

\subsection{Observaciones durante el entrenamiento}

Durante el entrenamiento se observó que:

\begin{itemize}
    \item La regresión logística convergió sin problemas dentro del límite de iteraciones establecido.
    \item El Random Forest no emitió advertencias de ningún tipo.
    \item Ambos modelos produjeron probabilidades en el rango [0, 1] para todas las observaciones, sin valores anómalos.
\end{itemize}

El hecho de que los modelos se entrenaran sin incidencias no garantiza que su rendimiento sea bueno, pero sí indica que la estructura del dataset era compatible con los algoritmos utilizados y que no existían problemas técnicos que impidieran el aprendizaje.

\section{Consideraciones sobre reproducibilidad}

A lo largo de todo el desarrollo se adoptaron medidas para garantizar la reproducibilidad del trabajo:

\begin{itemize}
    \item \textbf{Semillas aleatorias fijas}: todas las operaciones con componente aleatorio (partición de datos, entrenamiento de modelos) utilizaron semillas fijas documentadas.
    \item \textbf{Versionado de datos}: cada fase del pipeline generó un dataset intermedio almacenado como fichero independiente, permitiendo reconstruir cualquier estado anterior.
    \item \textbf{Documentación en notebooks}: cada decisión metodológica quedó documentada junto al código que la implementaba.
    \item \textbf{Serialización de modelos}: los modelos entrenados se guardaron en formato serializado para su reutilización exacta.
    \item \textbf{Control de versiones}: el proyecto se gestionó mediante Git, manteniendo un historial de cambios.
\end{itemize}

Estas prácticas permiten que cualquier persona con acceso al código y los datos originales pueda reproducir exactamente los resultados obtenidos, requisito fundamental en trabajos de investigación aplicada.

%% Ejemplo de figura
%% \begin{figure}[H]
%%     \centering
%%     \includegraphics[width=0.8\textwidth]{img/figura_ejemplo.png}
%%     \caption{Descripción de la figura. Fuente: elaboración propia.}
%%     \label{fig:ejemplo}
%% \end{figure}

%% Ejemplo de tabla
%% \begin{table}[H]
%%     \centering
%%     \caption{Descripción de la tabla. Fuente: elaboración propia.}
%%     \label{tab:ejemplo}
%%     \begin{tabular}{lcc}
%%         \toprule
%%         \textbf{Variable} & \textbf{Valor 1} & \textbf{Valor 2} \\
%%         \midrule
%%         Variable A & 0.85 & 0.90 \\
%%         Variable B & 0.78 & 0.82 \\
%%         \bottomrule
%%     \end{tabular}
%% \end{table}

%% --- CAPÍTULO 7: RESULTADOS ---
\chapter{Resultados}
\label{chap:resultados}

\section{Análisis descriptivo de la muestra}
\label{sec:descriptivo}

Características de la población de estudio.

\section{Rendimiento de los modelos}
\label{sec:rendimiento}

Comparación del rendimiento de los diferentes modelos.

\section{Interpretabilidad del modelo}
\label{sec:interpretabilidad}

Análisis de la importancia de variables y explicabilidad.

\section{Discusión}
\label{sec:discusion}

Análisis crítico de los resultados obtenidos y comparación con la literatura.

%% --- CAPÍTULO 8: CONCLUSIONES ---
\chapter{Conclusiones}
\label{chap:conclusiones}

\textit{Verificar si se han alcanzado los objetivos planteados al principio de esta memoria. En caso positivo, comentar las evidencias más relevantes y en caso negativo, indicar los motivos por los que se cree que ha fallado el objetivo en cuestión.}

\section{Conclusiones generales}
\label{sec:conclusiones_generales}

\textit{Conclusiones principales del trabajo.}

\section{Limitaciones del estudio}
\label{sec:limitaciones}

\textit{Limitaciones identificadas en el desarrollo del trabajo.}

%% --- CAPÍTULO 9: TRABAJOS FUTUROS ---
\chapter{Trabajos Futuros}
\label{chap:trabajos_futuros}

\textit{Describir cómo se le daría continuidad a este trabajo en el futuro, posibles nuevas líneas de investigación, o si el trabajo y metodología podrían extrapolarse a otros sectores o proyectos relacionados.}

%% ==================== BIBLIOGRAFÍA ==================== %%
\newpage
\printbibliography[heading=bibintoc, title={Referencias Bibliográficas}]

%% ==================== ANEXOS ==================== %%
\appendix
\chapter{Anexo A: Código fuente}
\label{chap:anexo_codigo}

\textit{Incluir aquí código fuente relevante o referencias al repositorio.}

\chapter{Anexo B: Tablas complementarias}
\label{chap:anexo_tablas}

\textit{Incluir aquí tablas adicionales que no se han incluido en el cuerpo principal.}

\end{document}
